% !TeX spellcheck = pt_BR
\documentclass[12pt,portuguese]{article}

%----------------------------------------------------------------------
% PACOTES ESSENCIAIS E TRADUÇÃO
%----------------------------------------------------------------------
\usepackage[brazilian]{babel}     % Suporte à língua portuguesa
\usepackage[utf8]{inputenc}       % Codificação de entrada
\usepackage[T1]{fontenc}          % Codificação da fonte
\usepackage{lmodern}              % Fonte escalável

%----------------------------------------------------------------------
% PACOTES MATEMÁTICOS
%----------------------------------------------------------------------
\usepackage{amsmath, amsthm, amssymb, amstext, mathtools}
\usepackage{mathrsfs}             % Para fontes como \mathscr
\usepackage{dsfont}               % Para o indicador de conjunto \mathds{1}
\usepackage[all]{xy}              % Para diagramas comutativos
\usepackage{imakeidx}             % Criação de Lista de Símbolos / Índice

% Desigualdades sempre com símbolos inclinados
\let\leq\leqslant
\let\le\leqslant
\let\ge\geqslant
\let\geq\geqslant

%----------------------------------------------------------------------
% FORMATAÇÃO, LAYOUT E CORES
%----------------------------------------------------------------------
\usepackage{geometry}
\usepackage{varwidth}
\usepackage{graphicx}
\usepackage[pdftex]{xcolor}
\usepackage{titlesec}
\usepackage{appendix}
\usepackage{enumitem}
\usepackage{float}
\usepackage{booktabs}
\usepackage{array}
\usepackage{colortbl}
\usepackage{fontawesome5}
\usepackage{lucide-icons}
\usepackage{caption}
\usepackage{stackengine}
\captionsetup{width=0.8\textwidth}

% Definição de Cores
\definecolor{teal}{RGB}{15,118,110}
\definecolor{text}{RGB}{28,25,23}
\definecolor{muted}{RGB}{87,83,78}
\definecolor{card}{RGB}{255,255,255}
\definecolor{border}{RGB}{231,229,228}

% Definição de margens e espaçamento
\setlength{\topmargin}{-0.7in}
\setlength{\oddsidemargin}{0in}
\setlength{\textheight}{9.55in}
\setlength{\textwidth}{6.5in}
\setlength{\baselineskip}{12mm}

% Formatação dos títulos de seção
\titleformat{\section}{\Large\bfseries}{\thesection.}{0.5em}{}
\titleformat{\subsection}{\large\bfseries}{\thesubsection.}{0.5em}{}
\titleformat{\subsubsection}{\normalfont\itshape}{\thesubsubsection.}{0.6em}{}

%----------------------------------------------------------------------
% TIKZ E GRÁFICOS
%----------------------------------------------------------------------
\usepackage{tikz}
\usepackage{pgfplots}
\pgfplotsset{compat=1.17}
\usetikzlibrary{positioning,calc,fit,nfold,decorations.pathreplacing,calligraphy,tikzmark,patterns,arrows.meta,bending,hobby}

%----------------------------------------------------------------------
% REFERÊNCIAS, LINKS E METADADOS DO PDF
%----------------------------------------------------------------------
\usepackage[pdftex, colorlinks, linkcolor={red!50!black}, 
citecolor={red}, urlcolor={blue!80!black}]{hyperref}

\hypersetup{
pdftitle={A Lei Forte dos Grandes Números},
pdfauthor={L. Cioletti},
pdfsubject={Lei forte dos grandes números, prova de Etemadi, truncamento, números normais},
pdfkeywords={lei forte dos grandes números, prova de Etemadi, truncamento, lema de Borel-Cantelli, números normais, MAT-UnB}
}

% --- PACOTE E COMANDO PARA A LICENÇA CREATIVE COMMONS ---
\usepackage[type={CC},modifier={by-nc-sa},version={4.0}]{doclicense}
\newcommand\blfootnote[1]{%
\begingroup
\renewcommand\thefootnote{}\footnote{#1}%
\addtocounter{footnote}{-1}%
\endgroup
}

%----------------------------------------------------------------------
% ALIASCNT E AMBIENTES DE TEOREMA (Cleveref)
%----------------------------------------------------------------------
\usepackage{aliascnt}
\theoremstyle{definition}
\newtheorem{theorem}{Teorema}

\newaliascnt{definition}{theorem}
\newtheorem{definition}[definition]{Definição}
\aliascntresetthe{definition}

\newaliascnt{example}{theorem}
\newtheorem{example}[example]{Exemplo}
\aliascntresetthe{example}

\newaliascnt{proposition}{theorem}
\newtheorem{proposition}[proposition]{Proposição}
\aliascntresetthe{proposition}

\newaliascnt{corollary}{theorem}
\newtheorem{corollary}[corollary]{Corolário}
\aliascntresetthe{corollary}

\newaliascnt{lemma}{theorem}
\newtheorem{lemma}[lemma]{Lema}
\aliascntresetthe{lemma}

\newaliascnt{remark}{theorem}
\newtheorem{remark}[remark]{Observação}
\aliascntresetthe{remark}

\newtheoremstyle{plainexercise}
  {3pt}{3pt}{\normalfont}{}{\bfseries}{.}{.5em}
  {\thmnumber{#2}\thmnote{ #3}}
\theoremstyle{plainexercise}
\newtheorem{exercise}{Exercício}

% Versão marcada dos exercícios: a estrela fica pendurada na margem esquerda,
% como nas notas anteriores.
\newtheoremstyle{starredexercise}
  {3pt}{3pt}{\normalfont}{}{\bfseries}{.}{.5em}
  {\llap{\textasteriskcentered\,}\thmnumber{#2}\thmnote{ #3}}
\theoremstyle{starredexercise}
\newtheorem{starredexercise}[exercise]{Exercício}
\theoremstyle{definition}

% Ambientes não numerados usados no texto
\newtheorem*{property}{Propriedade}
\newtheorem*{solution}{Solução}

% Se um enumerate for o primeiro conteúdo de um teorema/exercício (cabeçalho
% ainda pendente, \@inlabel true), o \leavevmode dispara o \everypar do amsthm
% que coloca o cabeçalho na própria linha, antes da lista iniciar.
\makeatletter
\setlist[enumerate]{before={\if@inlabel\leavevmode\fi}}
\makeatother

% Nomes em português para os ambientes
\AtBeginDocument{\renewcommand{\proofname}{\normalfont\textbf{Prova}}}
\renewcommand{\qedsymbol}{$\blacksquare$}

% Cleveref (deve vir depois de hyperref e amsthm)
\usepackage[brazilian,nameinlink]{cleveref}

\crefname{theorem}{teorema}{teoremas}
\Crefname{theorem}{Teorema}{Teoremas}
\crefname{definition}{definição}{definições}
\Crefname{definition}{Definição}{Definições}
\crefname{lemma}{lema}{lemas}
\Crefname{lemma}{Lema}{Lemas}
\crefname{proposition}{proposição}{proposições}
\Crefname{proposition}{Proposição}{Proposições}
\crefname{corollary}{corolário}{corolários}
\Crefname{corollary}{Corolário}{Corolários}
\crefname{example}{exemplo}{exemplos}
\Crefname{example}{Exemplo}{Exemplos}
\crefname{exercise}{exercício}{exercícios}
\Crefname{exercise}{Exercício}{Exercícios}
\crefname{remark}{observação}{observações}
\Crefname{remark}{Observação}{Observações}
\crefname{section}{seção}{seções}
\Crefname{section}{Seção}{Seções}

%----------------------------------------------------------------------
% MACROS ADICIONAIS
%----------------------------------------------------------------------
\let\oldemptyset\emptyset
\let\emptyset\varnothing
\graphicspath{{./}{../arquivos-aux/}}

% Seno na convenção em português
\renewcommand{\sin}{\operatorname{sen}}

%----------------------------------------------------------------------
% DOCUMENTO PRINCIPAL
%----------------------------------------------------------------------
\begin{document}
% --- CABEÇALHO COM O LOGO DA UnB ---
\vspace*{-0.8cm}
\begin{figure}[H]
\includegraphics{logo-mat-unb.pdf}
\end{figure}


% --- TÍTULO, AUTOR E DATA ---
\vspace*{0.3cm}
\begin{center}
{\LARGE\bf A Lei Forte dos Grandes Números}
\vspace{1cm}

{\large \bf L. Cioletti}
\vspace{0.3cm}

{\footnotesize Outubro de 2026}
\end{center}

% --- RESUMO (ABSTRACT) ---
\begin{abstract}
Nestas notas, estudamos a Lei Forte dos Grandes Números e suas aplicações.
Apresentamos inicialmente uma demonstração para variáveis aleatórias
independentes e identicamente distribuídas com quarto momento finito,
baseada na desigualdade de Markov e no critério de Borel-Cantelli. Em
seguida, demonstramos a versão de Etemadi para variáveis integráveis,
identicamente distribuídas e independentes dois a dois. Mostramos também
que, no caso independente e identicamente distribuído, a integrabilidade
é necessária para a convergência quase certa das médias a um limite finito.
Analisamos o comportamento das médias amostrais de variáveis independentes
com distribuição de Cauchy. A família de Pareto exemplifica a aplicação
da Lei Forte a sequências de variáveis integráveis com variância infinita.
Entre as aplicações, incluímos a normalidade dos números reais, a integração de Monte Carlo e a
convergência da função de distribuição empírica. As notas contêm exemplos
e exercícios propostos. A exposição segue as Seções 2.3 e 2.4 de
\cite{MR3930614}, as Seções 1, 2 e 6 do Capítulo 1 de \cite{MR2893652}
e as Seções 8.1 e 8.2 do Capítulo 8 de \cite{GrimmettWelsh2014}.
\end{abstract}

% --- LICENÇA NO RODAPÉ ---
\blfootnote{%
\doclicenseImage[imagewidth=2.6em]\ 
\begingroup
\hypersetup{urlcolor=darkgray}
\href{\doclicenseURL}{%
\doclicenseName - Attribution-NonCommercial-ShareAlike 4.0 International%
}%
\endgroup
}

%----------------------------------------------------------------------
% CONTEÚDO
%----------------------------------------------------------------------
\section{Introdução}\label{sec:introducao}

Nas notas anteriores sobre modos de convergência, aprendemos a distinguir
cinco noções de convergência para sequências de variáveis aleatórias. A lei
fraca dos grandes números, estudada nas notas anteriores afirma que, sob determinadas condições, 
a média empírica $S_n/n$ converge em probabilidade para seu valor esperado $\mu$, isto é, 
\begin{align*}
\frac{S_n}{n} \xrightarrow[\ n\to\infty\ ]{\mathbb{P}} \mu.
\end{align*}
Ou seja, para $n\in\mathbb{N}$ suficientemente grande, a probabilidade de que a
média se afaste de $\mu$ por uma magnitude maior que $\varepsilon$ é pequena. 
A lei forte faz uma afirmação ainda mais precisa, ela garante que para quase todo ponto $\omega\in\Omega$, 
a sequência numérica $S_n(\omega)/n$ converge para $\mu$ no sentido usual de sequências numéricas. 

A diferença entre as duas leis não é um refinamento técnico. Considere o
lançamento repetido de um dado honesto e seja $S_n$ a soma dos $n$ primeiros
resultados. A lei fraca diz que, para $n$ grande, é improvável que $S_n/n$ se
afaste de $7/2$. A lei forte diz que, com probabilidade $1$, existe um índice
a partir do qual $S_n(\omega)/n$ permanece próximo de $7/2$ para sempre. A
segunda afirmação implica a primeira, mas não é consequência dela: há
sequências que convergem em probabilidade e, ainda assim, oscilam infinitas
vezes em quase todos os pontos, como a sequência que varre o intervalo
discutida nas notas anteriores. A lei forte exclui esse comportamento para as
médias de variáveis independentes e identicamente distribuídas com
$\mathbb{E}[|X_1|]<\infty$.

O exemplo mais usado nestas notas para ilustrar os resultados principais
será o das frequências de dígitos. Mais precisamente, suponha que
$\omega\in [0,1]$ é um número real sorteado ao acaso, de maneira uniforme, e considere sua
expansão decimal $\omega=0{,}d_1d_2d_3\ldots$. Qual é a frequência relativa
do dígito $7$ entre os $n$ primeiros dígitos, quando $n$ cresce? A intuição diz
que a frequência deve se aproximar de $1/10$, e o teorema de Borel, que
demonstraremos, confirma a intuição para quase todo $\omega$. O mesmo problema
pode ser formulado de maneira mais clara usando o espaço produto $(\Omega,\mathcal{F},\mathbb{P})$, onde
\[
\Omega\equiv \{0,1,\ldots,9\}^{\mathbb{N}},
\]
$\mathcal{F}$ é uma $\sigma$-álgebra apropriada e a medida de probabilidade $\mathbb{P}$
sendo a única medida de probabilidade sobre $\mathcal{F}$ 
satisfazendo $\mathbb{P}(X_n=k)=1/10$, onde
$X_n(\omega)=\omega_n$ é a $n$-ésima projeção $X_n:\Omega\to \{0,1,2,\ldots,9\}$ 
(que sobre o espaço $(\Omega,\mathcal{F},\mathbb{P})$ é uma variável aleatória)
e $k\in \{0,1,2,\ldots,9\}$ é um dígito fixado. 
Neste modelo, os dígitos são representados pelas variáveis aleatórias $X_n$'s 
que são, por construção, independentes e identicamente distribuídas. 
Além do mais, neste modelo a frequência relativa de um dígito $k$ fixado, 
é simplesmente a média empírica da sequência de indicadoras $\mathds{1}_{\{X_j=k\}}$,
isto é, 
\begin{align*}
\frac{1}{n} \sum_{j=1}^n \mathds{1}_{\{X_j=k\}}
\end{align*}
Neste modelo o problema da determinação da frequência relativa dos dígitos
pode ser obtido como uma aplicação da Lei Forte dos Grandes Números. 


A estratégia da prova segue de perto a adotada na referência \cite{MR3930614}. 
Primeiro, apresentamos uma belíssima prova da Lei Forte para o caso em que as
variáveis têm quarto momento finito. Além da prova neste caso ser muito clara
e curta ela ilustra como as principais ferramentas são utilizadas no caso geral. 
Veremos que se assumimos a existência do quarto momento,
então a desigualdade de Markov aplicada à quarta potência da soma centrada, 
$(S_n-n\mu)^4$, fornece uma cota de ordem $1/n^2$, cuja soma converge. 
Este fato irá permitir aplicar o critério de Borel-Cantelli, 
apresentado nas notas anteriores, para concluir a ocorrência da convergência quase
certa.
Embora esta versão seja ainda um caso particular ela é suficiente 
para resolver o problema da frequência relativa do dígitos, porque as variáveis
indicadoras são limitadas e portanto possuem todos os momentos finitos. 

Depois, voltamos nossas atenções para o caso geral. 
Vamos então substituir a hipótese de quarto momento por 
$\mathbb{E}[|X_1|]<\infty$. A prova apresentada nestas notas 
para este caso é a prova inspirada no argumento de Etemadi. Como veremos, este argumento
contorna a ausência de momentos de ordem superior por meio de um truncamento e de uma
estimativa de momentos para as variáveis truncadas. O truncamento faz com que
cada variável seja limitada por uma escala que cresce com o índice, e a
estimativa mostra que a soma das variâncias truncadas, dividida pelos
quadrados dos índices, é finita. A convergência das esperanças truncadas
identifica o centro dessas somas com a média original. Com isso, a cota de
Chebyshev tem somabilidade garantida ao longo de uma subsequência geométrica
de índices, mesmo quando não a temos ao longo de todos os índices. O
preenchimento dos índices restantes usa a monotonicidade da soma truncada, e é
aqui que a redução a variáveis não negativas se torna necessária.

A prova de Etemadi é um bom exemplo de argumento em que cada etapa tem uma
função clara: o truncamento elimina a cauda das variáveis, a estimativa de
momentos controla a variância acumulada, a subsequência geométrica torna a
série de probabilidades convergente e a monotonicidade transporta a conclusão
da subsequência para todos os índices. Ao longo da exposição, procuraremos
explicitar qual dessas funções cada parte do argumento desempenha, porque a prova é
frequentemente apresentada como uma sequência de estimativas sem que o leitor
veja por que cada passo foi dado.

O texto está organizado da seguinte maneira.
Na \Cref{sec:preliminares} recordamos
as ferramentas de cálculo de probabilidades que serão usadas: as desigualdades
de Markov e Chebyshev, os dois lemas de Borel-Cantelli, o critério de
convergência quase certa por somabilidade e a representação da esperança pela
integral da cauda. Na seção \Cref{sec:espaco-produto} apresentamos mais alguns detalhes 
sobre a construção do espaço produto e sobre o problema das frequências relativas. 
Na \Cref{sec:quarto-momento} demonstramos a lei forte sob a hipótese de quarto momento finito 
e em seguida, mostramos como este resultado é aplicado para resolver o problema
da frequência relativa dos dígitos. 
Na \Cref{sec:etemadi} apresentamos a prova de Etemadi.
A \Cref{sec:reciproca} mostra que a hipótese $\mathbb{E}[|X_1|]<\infty$ é
necessária no caso independente e identicamente distribuído.
A \Cref{sec:aplicacoes} reúne aplicações, e a \Cref{sec:exercicios} propõe exercícios.

A exposição da prova de Etemadi segue de perto a Seção 2.4 de
\cite{MR3930614}, com detalhamento maior das passagens que o livro deixa
para o leitor. O exemplo dos dígitos e o enunciado do Teorema de Borel seguem as
Seções 1 e 2 do Capítulo 1 de \cite{MR2893652}, adaptadas ao espaço produto
$\{0,1,\ldots,9\}^{\mathbb{N}}$. A prova da Lei Forte no caso de quarto momento é 
inspirada na Seção 2.3 da referência \cite{MR3930614}.

%----------------------------------------------------------------------
\section{Preliminares}\label{sec:preliminares}

Nesta seção, reunimos as ferramentas que serão usadas nas provas da lei forte
e de sua recíproca. Recordamos as desigualdades e os lemas das notas anteriores
e explicitamos a relação entre integrabilidade e probabilidades de cauda.

\subsection{Desigualdades de Markov e Chebyshev e a Lei Fraca}\label{subsec:markov-chebyshev}

As desigualdades de Markov e Chebyshev, estudadas em notas de aula anteriores,
constituem o ponto de partida natural para estimar probabilidades de cauda e
desvios a partir de momentos. A pergunta fundamental que orienta essas
estimativas é simples: conhecendo apenas a esperança ou a variância de uma
variável aleatória, o que podemos garantir sobre a probabilidade de ela assumir
valores excessivamente grandes ou afastar-se de sua média?

A resposta mais elementar é dada pela desigualdade de Markov: se $Z$ é uma
variável aleatória não negativa e $a>0$, então
\begin{equation}\label{eq:markov}
\mathbb{P}(Z\geqslant a)\leqslant\frac{\mathbb{E}[Z]}{a}.
\end{equation}
A intuição por trás de \eqref{eq:markov} é imediata: se a probabilidade de $Z$
superar um patamar $a$ fosse estritamente maior que $\mathbb{E}[Z]/a$, a
contribuição desse evento já excederia a própria média global de $Z$, o que é
impossível.

Quando dispomos de mais informação sobre a distribuição, como o segundo momento
finito, podemos refinar substancialmente o controle das flutuações. Seja $X$
uma variável aleatória com $\mathbb{E}[X^2]<\infty$. Fixado $\varepsilon>0$,
a função quadrática nos permite transformar o evento de desvio
$\{|X-\mathbb{E}[X]|\geqslant\varepsilon\}$ na condição equivalente
$\{(X-\mathbb{E}[X])^2\geqslant\varepsilon^2\}$. Aplicando a desigualdade de
Markov \eqref{eq:markov} à variável não negativa $Z=(X-\mathbb{E}[X])^2$ com o
patamar $a=\varepsilon^2$, obtemos a desigualdade de Chebyshev:
\begin{equation}\label{eq:chebyshev}
\mathbb{P}\bigl(|X-\mathbb{E}[X]|\geqslant\varepsilon\bigr)
\leqslant
\frac{\operatorname{Var}(X)}{\varepsilon^2}.
\end{equation}
A desigualdade \eqref{eq:chebyshev} mostra com clareza como a variância mede a
concentração de $X$ em torno de sua esperança: quanto menor a dispersão, menor
a probabilidade de observarmos desvios de ordem $\varepsilon$.

Vejamos agora como essa estimativa se comporta ao avaliarmos o comportamento das
médias amostrais. Considere uma sequência de variáveis aleatórias $X_1,X_2,\ldots$
com variâncias finitas e suponha que sejam duas a duas independentes. Definindo a
soma parcial $S_n=X_1+\cdots+X_n$, a independência aos pares assegura que as
covariâncias se anulam, de modo que a variância da soma se reduz à soma das
variâncias:
\[
\operatorname{Var}(S_n)=\sum_{k=1}^{n}\operatorname{Var}(X_k).
\]
Consequentemente, a variância da média amostral $S_n/n$ decai com o quadrado de
$n$ no denominador:
\[
\operatorname{Var}\Bigl(\frac{S_n}{n}\Bigr)
=
\frac{1}{n^2}\sum_{k=1}^{n}\operatorname{Var}(X_k).
\]
Aplicando a desigualdade de Chebyshev \eqref{eq:chebyshev} à variável $S_n/n$,
obtemos, para cada $\varepsilon>0$, a estimativa
\begin{equation}\label{eq:lei-fraca}
\mathbb{P}\Bigl(\Bigl|\frac{S_n}{n}-\bar\mu_n\Bigr|\geqslant\varepsilon\Bigr)
\leqslant
\frac{1}{\varepsilon^2}\operatorname{Var}\Bigl(\frac{S_n}{n}\Bigr),
\qquad\text{onde}\qquad
\bar\mu_n=\frac{1}{n}\sum_{k=1}^{n}\mathbb{E}[X_k].
\end{equation}
No cenário particular em que as variáveis são identicamente distribuídas, com
média comum $\mu$ e variância finita $\sigma^2$, temos $\bar\mu_n=\mu$ e o lado
direito de \eqref{eq:lei-fraca} assume a forma explícita
\[
\mathbb{P}\Bigl(\Bigl|\frac{S_n}{n}-\mu\Bigr|\geqslant\varepsilon\Bigr)
\leqslant
\frac{\sigma^2}{n\varepsilon^2}.
\]
Como essa cota tende a zero quando $n\to\infty$, concluímos que a média amostral
$S_n/n$ converge em probabilidade para a média teórica $\mu$. Esse resultado é
precisamente a Lei Fraca dos Grandes Números sob a hipótese de variância finita.

No entanto, há uma limitação crucial nessa estimativa que motiva todo o estudo
que faremos a seguir. Sempre que $\sigma^2>0$, a cota superior fornecida por
Chebyshev decai apenas com taxa $1/n$. Como a série harmônica $\sum_{n=1}^\infty 1/n$
diverge, essas probabilidades de desvio não são somáveis em $n$. Isso nos impede
de usar diretamente o primeiro lema de Borel-Cantelli para deduzir que a
convergência ocorre quase certamente. Essa constatação é a motivação central
destas notas de aula: para demonstrar a Lei Forte dos Grandes Números, precisaremos
ou de momentos superiores que produzam cotas somáveis (como $1/n^2$), ou de
estratégias mais refinadas que explorem truncamentos e subsequências ao longo
das quais a somabilidade possa ser recuperada.

\subsection{Borel-Cantelli e o Critério de Convergência Quase Certa}
\label{subsec:bc-criterio}

Os resultados a seguir foram discutidos nas notas \emph{Modos de Convergência
de Sequências de Variáveis Aleatórias}. Aqui recordamos os enunciados na forma
em que serão usados nas provas da lei forte e de sua recíproca.

Para uma sequência de eventos $(A_n)$, o limite superior
\[
\limsup_{n\to\infty}A_n
=\bigcap_{N=1}^{\infty}\bigcup_{n\geqslant N}A_n
\]
é o evento formado pelos pontos $\omega$ que pertencem a $A_n$ para infinitos
índices $n$. Escrevemos também $\{A_n\ \text{i.v.}\}$ para esse
evento. O primeiro lema de Borel-Cantelli afirma que, se
$\sum_n\mathbb{P}(A_n)<\infty$, então
\[
\mathbb{P}\Bigl(\limsup_{n\to\infty}A_n\Bigr)=0,
\]
isto é, com probabilidade um apenas finitos dos eventos $A_n$ ocorrem. 
Note que este resultado não exige a independência dos eventos.

\medskip 

Aplicado aos eventos de desvio $\{|X_n-X|\geqslant\varepsilon\}$, o 
primeiro lema fornece o \emph{critério de Borel-Cantelli para a convergência quase
certa}: se $X,X_1,X_2,\ldots$ são variáveis aleatórias e, para todo
$\varepsilon>0$,
\begin{equation}\label{eq:criterio}
\sum_{n=1}^{\infty}\mathbb{P}\bigl(|X_n-X|\geqslant\varepsilon\bigr)<\infty,
\end{equation}
então $X_n\to X$ quase certamente. A demonstração detalhada desse critério
foi apresentada na Proposição 8 da Subseção~3.5 das notas de aula
\href{https://mat.unb.br/~cioletti/ensino/modos-convergencia.pdf}{\emph{Modos de Convergência
de Sequências de Variáveis Aleatórias}}.

A convergência em probabilidade pede apenas que a probabilidade de cada desvio
tenda a zero; a somabilidade em \eqref{eq:criterio} permite controlar sua
repetição ao longo da sequência. Essa condição é suficiente para a
convergência quase certa, mas não necessária, como mostram os exemplos da
Subseção~3.6, \emph{O Critério não é Necessário}, das notas anteriores. Nas duas
provas da lei forte que apresentaremos, procuraremos estimativas que permitam
aplicar o critério.

Na prova da recíproca, usaremos também o segundo lema de Borel-Cantelli: se os
eventos $A_n$ são independentes e $\sum_n\mathbb{P}(A_n)=\infty$, então
\[
\mathbb{P}\Bigl(\limsup_{n\to\infty}A_n\Bigr)=1.
\]
Nesse caso, quase certamente infinitos dos eventos $A_n$ ocorrem. Os dois
lemas estão enunciados na Seção~2 das notas anteriores.

\subsection{Esperança e Probabilidades de Cauda}\label{subsec:esperanca-cauda}

Se $W$ é uma variável aleatória não negativa, a fórmula da esperança por
integração da cauda afirma que
\begin{equation}\label{eq:esperanca-cauda}
\mathbb{E}[W]=\int_0^{\infty}\mathbb{P}(W>t)\,dt,
\end{equation}
com ambos os lados podendo ser infinitos. De fato,
\[
W(\omega)=\int_0^{\infty}\mathds{1}_{(0,W(\omega))}(t)\,dt,
\]
e o teorema de Tonelli permite trocar a esperança com a integral, pois o
integrando é não negativo.

Como $t\mapsto\mathbb{P}(W>t)$ é não crescente, para todo inteiro $k\geqslant0$
temos
\[
\mathbb{P}(W>k+1)
\leqslant\int_k^{k+1}\mathbb{P}(W>t)\,dt
\leqslant\mathbb{P}(W>k).
\]
Somando essas desigualdades e usando \eqref{eq:esperanca-cauda}, obtemos
\begin{equation}\label{eq:comparacao-cauda}
\sum_{k=1}^{\infty}\mathbb{P}(W>k)
\leqslant\mathbb{E}[W]
\leqslant\sum_{k=0}^{\infty}\mathbb{P}(W>k).
\end{equation}
As séries dos extremos diferem apenas pelo termo $\mathbb{P}(W>0)\leqslant1$.
Portanto, $\mathbb{E}[W]<\infty$ se, e somente se,
$\sum_{k=1}^{\infty}\mathbb{P}(W>k)<\infty$. Aplicaremos essa equivalência a
variável aleatória
$W \equiv |X_1|$. A integrabilidade tornará raros os truncamentos na prova de
Etemadi, enquanto sua ausência produzirá termos grandes infinitas vezes na
prova da recíproca.

%----------------------------------------------------------------------
\section{O Espaço Produto e as Frequências dos Dígitos}\label{sec:espaco-produto}

\subsection{A Construção do Modelo}\label{subsec:modelo}

Trabalharemos no espaço de todas as sequências infinitas de dígitos,
\[
\Omega=\{0,1,\ldots,9\}^{\mathbb{N}}
=
\bigl\{\omega=(\omega_1,\omega_2,\omega_3,\ldots):
\omega_n\in\{0,1,\ldots,9\}\text{ para todo }n\bigr\}.
\]
As funções $X_n:\Omega\to \mathbb{R}$ dadas por 
\[
X_n(\omega)=\omega_n,
\qquad n\geqslant1,
\]
podem ser pensadas como variáveis aleatórias sobre o espaço mensurável $(\Omega,\mathcal{F})$, 
onde $\mathcal{F}$ é a menor $\sigma$-álgebra, onde para todo $n\in\mathbb{N}$, 
cada função $X_n$ é mensurável. 


Um \emph{cilindro de posto $n$} é um subconjunto de $\Omega$ obtido fixando as $n$ primeiras
coordenadas:
\[
C(u_1,\ldots,u_n)
=
\bigl\{\omega\in\Omega:X_1(\omega)=u_1,\ldots,X_n(\omega)=u_n\bigr\},
\]
com $u_1,\ldots,u_n\in\{0,1,\ldots,9\}$. É claro pela definição de $\mathcal{F}$
que qualquer cilindro da forma $C(u_1,\ldots,u_n)\in \mathcal{F}$. Basta
observar que este cilindro corresponde ao evento 
\begin{align*}
\{X_1=u_1\}\cap \ldots\cap \{X_n=u_n\}\in\mathcal{F}.
\end{align*}

A medida de probabilidade $\mathbb{P}$ que nos interessa sobre $(\Omega,\mathcal{F})$ 
é a única medida de probabilidade definida sobre o espaço mensurável $(\Omega,\mathcal{F})$ 
satisfazendo a seguinte propriedade
\begin{equation}\label{eq:medida-produto}
\mathbb{P}\bigl(C(u_1,\ldots,u_n)\bigr)=\frac{1}{10^n}
=
\prod_{j=1}^{n}\frac{1}{10}
\end{equation}

\begin{proposition}\label{prop:marginais}
Para todo $n\geqslant1$ e todo $k\in\{0,1,\ldots,9\}$,
\[
\mathbb{P}(X_n=k)=\frac{1}{10}.
\]
\end{proposition}

\begin{proof}
O evento $\{X_n=k\}$ é a união disjunta dos $10^{n-1}$ cilindros de posto $n$
cuja $n$-ésima coordenada vale $k$, um para cada escolha das $n-1$ primeiras
coordenadas. Por \eqref{eq:medida-produto}, cada um desses cilindros tem
probabilidade $10^{-n}$. Logo
\[
\mathbb{P}(X_n=k)=10^{n-1}\cdot10^{-n}=\frac{1}{10}.
\]
\end{proof}

\begin{proposition}\label{prop:independencia}
A sequência $(X_n)_{n\geqslant1}$ é independente. Mais precisamente, para todos
$1\leqslant n_1<n_2<\cdots<n_j$ e todos
$k_1,\ldots,k_j\in\{0,1,\ldots,9\}$,
\[
\mathbb{P}\bigl(X_{n_1}=k_1,\ldots,X_{n_j}=k_j\bigr)
=
\prod_{i=1}^{j}\mathbb{P}(X_{n_i}=k_i).
\]
\end{proposition}

\begin{proof}
O evento $\{X_{n_1}=k_1,\ldots,X_{n_j}=k_j\}$ é a união disjunta dos cilindros
de posto $n_j$ compatíveis com as $j$ condições. As coordenadas fora do conjunto
$\{n_1,\ldots,n_j\}$ são livres, em número de $n_j-j$, de modo que há
$10^{\,n_j-j}$ cilindros. Por \eqref{eq:medida-produto},
\[
\mathbb{P}\bigl(X_{n_1}=k_1,\ldots,X_{n_j}=k_j\bigr)
=
\frac{10^{\,n_j-j}}{10^{\,n_j}}
=
10^{-j}
=
\prod_{i=1}^{j}\frac{1}{10},
\]
e a última expressão é o produto das probabilidades marginais pela
\Cref{prop:marginais}.
\end{proof}

Cada evento $\{X_n\in B\}$, com $B\subseteq\{0,1,\ldots,9\}$, é uma união
finita disjunta de eventos $\{X_n=k\}$. Somando a fatoração demonstrada acima,
obtemos a independência para subconjuntos arbitrários e, portanto,
para toda a sequência. O modelo é o de uma sequência de sorteios independentes
de um dígito, cada um com probabilidade $1/10$. Ele é a versão de dez alternativas do
modelo de lançamentos de uma moeda honesta e também pode ser obtido a partir da
expansão decimal de um ponto escolhido ao acaso em $[0,1]$, como veremos na
\Cref{subsec:digitos-reais}.

\subsection{O Problema das Frequências}\label{subsec:problema-frequencias}

Para $n\geqslant1$ e $k\in\{0,1,\ldots,9\}$, definimos a frequência relativa do
dígito $k$ entre os $n$ primeiros dígitos:
\begin{equation}\label{eq:frequencia}
\nu_{n,k}(\omega)
=
\frac{1}{n}\sum_{j=1}^{n}\mathds{1}_{\{X_j=k\}}(\omega).
\end{equation}
A variável $\nu_{n,k}$ é a proporção das $n$ primeiras posições ocupadas pelo
dígito $k$. Pela \Cref{prop:marginais}, cada indicadora tem esperança
$1/10$, logo
\[
\mathbb{E}[\nu_{n,k}]=\frac{1}{10}
\qquad\text{para todo }n.
\]
Resta saber se a frequência relativa converge para esse valor ao longo da
sequência, e em que sentido.

\begin{theorem}[Teorema de Borel para a base 10]\label{thm:borel}
Para cada $k\in\{0,1,\ldots,9\}$,
\[
\lim_{n\to\infty}\nu_{n,k}(\omega)=\frac{1}{10}
\qquad\text{para quase todo }\omega\in\Omega.
\]
Além disso, a interseção dos conjuntos de convergência sobre os dez valores de
$k$ tem probabilidade $1$, de modo que, quase certamente, todos os dez dígitos
têm frequência assintótica $1/10$ simultaneamente.
\end{theorem}

A segunda afirmação é consequência da primeira e da subaditividade da
probabilidade: a interseção de dez eventos de probabilidade $1$ tem
probabilidade $1$. Registramos o resultado como um único teorema porque é assim
que ele é usado.

\begin{definition}[Sequências simplesmente normais e normais]\label{def:normal}
Um ponto $\omega\in\Omega$ é dito \emph{simplesmente normal na base $10$} se,
para todo $k\in\{0,1,\ldots,9\}$,
\[
\lim_{n\to\infty}\nu_{n,k}(\omega)=\frac{1}{10}.
\]
Dizemos que $\omega$ é \emph{normal na base $10$} se, para toda palavra finita
$u=u_1\ldots u_\ell$ de dígitos com $\ell\geqslant1$, a frequência assintótica
de ocorrências do bloco $u$ nas posições consecutivas de $\omega$ é igual a
$10^{-\ell}$.
\end{definition}

Na definição de normalidade, todas as posições iniciais são levadas em conta,
permitindo ocorrências sobrepostas do bloco $u$. Enquanto a normalidade simples
garante apenas o equilíbrio assintótico de cada dígito isolado, a normalidade
plena impõe que qualquer sequência finita de dígitos ocorra exatamente com a
frequência ditada pela probabilidade uniforme.

Com essa terminologia, o \Cref{thm:borel} afirma precisamente que o conjunto dos
pontos simplesmente normais na base $10$ tem probabilidade um, ou seja, 
uma sequência sorteada ao acaso satisfaz essa propriedade de
equidistribuição.

O \Cref{thm:borel} decorre naturalmente da Lei Forte dos Grandes Números. Na
\Cref{sec:quarto-momento}, demonstramos a lei forte sob a hipótese de quarto
momento finito, que pode portanto ser aplicado ao caso de sequências de
variáveis aleatórias limitadas, como é o caso das indicadoras dos dígitos.
No entanto, a exigência de momentos de ordem quatro pode ser bastante relaxada. Na
\Cref{sec:etemadi}, apresentamos a demonstração da lei forte geral sob a
hipótese mínima de primeiro momento finito ($\mathbb{E}[|X_1|]<\infty$), devida a
Etemadi. Mais adiante, na \Cref{sec:aplicacoes}, exploramos a força dessa
generalidade ao estudar distribuições de cauda pesada, como a distribuição de
Pareto com expoente em $(1,2]$, cujas médias amostrais convergem quase
certamente embora todos os momentos de ordem superior sejam infinitos.

Antes de prosseguir, vale registrar uma sutileza do modelo. A expansão decimal
de um número real não é única: $0{,}4999\ldots=0{,}5000\ldots$. O conjunto dos
pontos com duas expansões é enumerável, logo tem probabilidade zero para a
medida uniforme em $[0,1]$, e a ambiguidade não afeta nenhuma afirmação quase
certa. No espaço produto não há ambiguidade alguma, pois cada $\omega$ é uma
sequência e não um número. Essa é uma das razões para trabalhar no espaço
produto: ele elimina uma complicação que é irrelevante para o problema.

%----------------------------------------------------------------------
\section{A Lei Forte com Quarto Momento}\label{sec:quarto-momento}

\subsection{Enunciado e Prova do Teorema}\label{subsec:quarto-momento-teorema}

Iniciamos nosso estudo da Lei Forte dos Grandes Números analisando o caso em que
as variáveis aleatórias possuem quarto momento finito. Essa demonstração é
particularmente instrutiva porque ilustra, em sua forma mais transparente e sem
as dificuldades técnicas de truncamento, como o critério de Borel-Cantelli
transforma estimativas de momentos em convergência quase certa.

Para compreender a estratégia que adotaremos, recordemos a dificuldade central
encontrada na \Cref{subsec:markov-chebyshev}: ao tentar estimar as probabilidades
de desvio pela desigualdade de Chebyshev (que se baseia no segundo momento e na
variância), obtivemos uma cota proporcional a $1/n$. Como a série harmônica
$\sum_{n=1}^\infty 1/n$ diverge, essa cota baseada no segundo momento não
permite aplicar diretamente o primeiro lema de Borel-Cantelli.

A ideia fundamental para contornar essa limitação consiste em explorar momentos
superiores. Aplicando a desigualdade de Markov à quarta potência do desvio,
$(S_n-n\mu)^4$, mostraremos que a esperança correspondente cresce apenas com
$n^2$. Dividindo por $n^4$, obteremos para a probabilidade de desvio uma cota
superior proporcional a $1/n^2$. Como a série $\sum_{n=1}^\infty 1/n^2$
converge, o critério de Borel-Cantelli garantirá de imediato que os desvios
ocorrem apenas um número finito de vezes, entregando a convergência quase certa.

\begin{theorem}[Lei forte com quarto momento]\label{thm:quarto-momento}
Sejam $X_1,X_2,\ldots$ variáveis aleatórias independentes e identicamente
distribuídas com $\mu=\mathbb{E}[X_i]$ e $\mathbb{E}[X_i^4]<\infty$. Se
$S_n=X_1+\cdots+X_n$, então
\[
\frac{S_n}{n}\longrightarrow\mu
\qquad\text{quase certamente.}
\]
\end{theorem}

\begin{proof}
Substituindo $X_i$ por $X_i-\mu$, podemos supor sem perda de generalidade que
$\mu=0$. De fato, se $X_i'=X_i-\mu$, então $S_n'=S_n-n\mu$ e a convergência
$S_n'/n\to0$ é equivalente a $S_n/n\to\mu$; além disso, as variáveis $X_i'$
continuam independentes, identicamente distribuídas e com quarto momento
finito.

Suponha então $\mu=0$. Expandindo o quarto momento da soma e usando a propriedade
de linearidade da esperança ficamos com 
\begin{equation}\label{eq:expansao-quarta}
\mathbb{E}[S_n^4]
=
\mathbb{E}\Bigl[\Bigl(\sum_{i=1}^{n}X_i\Bigr)^4\Bigr]
=
\sum_{i,j,k,\ell=1}^{n}\mathbb{E}[X_iX_jX_kX_\ell].
\end{equation}
Vamos classificar as parcelas de acordo com o número de índices distintos.

Se algum índice aparece uma única vez entre $i,j,k,\ell$, a parcela se anula.
Por exemplo, se $i,j,k$ são distintos e $\ell=i$, então
$\mathbb{E}[X_i^2X_jX_k]=\mathbb{E}[X_i^2]\,\mathbb{E}[X_j]\,\mathbb{E}[X_k]=0$,
porque $\mathbb{E}[X_j]=0$ e as variáveis são independentes. O mesmo argumento
se aplica a $\mathbb{E}[X_i^3X_j]$ e a $\mathbb{E}[X_iX_jX_kX_\ell]$ com os
quatro índices distintos.

Restam duas formas. A primeira é a dos termos com os quatro índices iguais,
$\mathbb{E}[X_i^4]$, que ocorrem $n$ vezes, uma para cada escolha de $i$. A
segunda é a dos termos com dois pares de índices iguais,
$\mathbb{E}[X_i^2X_j^2]=(\mathbb{E}[X_i^2])^2$ com $i\neq j$. Para contá-los,
escolhemos o par $\{i,j\}$ de $\binom{n}{2}$ maneiras e, uma vez fixado o par, o
termo $X_i^2X_j^2$ aparece em $6$ posições distintas, correspondentes às
permutações de $i,i,j,j$. Logo há
\[
6\binom{n}{2}=3n(n-1)
\]
tais termos. Como as variáveis são identicamente distribuídas,
\begin{equation}\label{eq:cota-quarta}
\mathbb{E}[S_n^4]
=
n\,\mathbb{E}[X_1^4]+3n(n-1)\bigl(\mathbb{E}[X_1^2]\bigr)^2
\leqslant
Cn^2,
\end{equation}
com
\[
C=\mathbb{E}[X_1^4]+3\bigl(\mathbb{E}[X_1^2]\bigr)^2<+\infty.
\]
Observe que $C$ não depende de $n$.

Fixe $\varepsilon>0$. Aplicamos agora a desigualdade de Markov à variável
não negativa $S_n^4$, com nível $\varepsilon^4n^4$:
\[
\mathbb{P}\Bigl(\Bigl|\frac{S_n}{n}\Bigr|\geqslant\varepsilon\Bigr)
=
\mathbb{P}\bigl(|S_n|\geqslant\varepsilon n\bigr)
=
\mathbb{P}\bigl(S_n^4\geqslant\varepsilon^4n^4\bigr)
\leqslant
\frac{\mathbb{E}[S_n^4]}{\varepsilon^4n^4}
\leqslant
\frac{C}{\varepsilon^4n^2}.
\]
Somando em $n$,
\begin{equation}\label{eq:serie-quarta}
\sum_{n=1}^{\infty}\mathbb{P}\Bigl(\Bigl|\frac{S_n}{n}-0\Bigr|\geqslant\varepsilon\Bigr)
\leqslant
\frac{C}{\varepsilon^4}\sum_{n=1}^{\infty}\frac{1}{n^2}
<
\infty.
\end{equation}
Como essa série converge para todo $\varepsilon>0$, segue diretamente do critério de
Borel-Cantelli (\Cref{subsec:bc-criterio}) que $S_n/n\to0$ quase certamente.
\end{proof}

\begin{remark}[Por que o quarto momento?]\label{rem:por-que-quarto-momento}
Poderíamos indagar se a mesma demonstração elementar funcionaria assumindo
apenas a existência de momentos de ordem inferior, como o terceiro momento
($\mathbb{E}[|X_1|^3]<\infty$). A resposta é negativa, e a estratégia empregada
quebra por duas razões estruturais:
\begin{enumerate}
\item[(i)] \textbf{Paridade da potência e não negatividade:} a função $x\mapsto x^4$
é par e não negativa, o que permitiu transformar diretamente o evento de desvio
$\{|S_n|\geqslant\varepsilon n\}$ em $\{S_n^4\geqslant\varepsilon^4 n^4\}$ e
aplicar a desigualdade de Markov à variável não negativa $S_n^4$. Com o terceiro
momento, a variável $S_n^3$ assume valores negativos, e para aplicar Markov
seríamos forçados a considerar a variável não negativa $|S_n|^3$.
\item[(ii)] \textbf{Álgebra da independência:} para calcular $\mathbb{E}[S_n^4]$,
expandimos a potência polinomialmente e usamos que a independência anula todas
as parcelas com índices ímpares isolados. Em contrapartida, a expressão
$|S_n|^3=|\sum_{i=1}^n X_i|^3$ não admite expansão em produtos algébricos simples
das variáveis $X_i$, destruindo o mecanismo combinatório elementar que zerou a
maioria dos termos cruzados.
\end{enumerate}
Portanto, o expoente $p=4$ é o menor inteiro par em que esse método algébrico
direto produz uma cota somável (de ordem $1/n^2$). Para $p=2$, a cota de
Chebyshev decai com taxa $1/n$ e não é somável. Já para contornar a exigência de
momentos superiores e alcançar a hipótese ótima $\mathbb{E}[|X_1|]<\infty$,
precisaremos de técnicas analíticas mais profundas, como o método de Etemadi
que apresentaremos na \Cref{sec:etemadi}.
\end{remark}

\begin{corollary}[Variáveis aleatórias limitadas]\label{cor:limitadas}
Se $(X_n)$ é uma sequência de variáveis aleatórias
independente e identicamente distribuída com $|X_n|\leqslant M$
quase certamente, e $\mu=\mathbb{E}[X_1]$, então
\begin{align*}
\frac{S_n}{n} \xrightarrow[\ n\to\infty \ ] \,  \mu, \qquad \text{quase certamente.}
\end{align*}
\end{corollary}

\begin{proof}
A hipótese de limitação implica imediatamente $\mathbb{E}[X_1^4]\leqslant M^4<\infty$.
Como as demais hipóteses do
\Cref{thm:quarto-momento} também são satisfeitas o resultado segue.
\end{proof}

\subsection{Aplicação às Frequências dos Dígitos}\label{subsec:digitos-quarta}

Agora combinamos o \Cref{thm:quarto-momento} com o \Cref{cor:limitadas} para
demonstrar o \Cref{thm:borel}. Fixe
$k\in\{0,1,\ldots,9\}$ e defina
\[
Y_j=\mathds{1}_{\{X_j=k\}}-\frac{1}{10},
\qquad j\geqslant1.
\]
Pela \Cref{prop:independencia}, as variáveis $Y_j$ são independentes e
identicamente distribuídas. Além disso,
\[
\mathbb{E}[Y_j]=\frac{1}{10}-\frac{1}{10}=0,
\qquad
|Y_j|\leqslant1,
\]
de modo que $\mathbb{E}[Y_j^4]\leqslant1<\infty$. Pelo
\Cref{cor:limitadas},
\begin{equation}\label{eq:digitos-media}
\frac{1}{n}\sum_{j=1}^{n}Y_j
\longrightarrow0
\qquad\text{quase certamente.}
\end{equation}
Mas
\[
\frac{1}{n}\sum_{j=1}^{n}Y_j
=
\frac{1}{n}\sum_{j=1}^{n}\mathds{1}_{\{X_j=k\}}-\frac{1}{10}
=
\nu_{n,k}-\frac{1}{10},
\]
logo \eqref{eq:digitos-media} é exatamente a afirmação
$\nu_{n,k}\to1/10$ quase certamente. Como $k$ é arbitrário e a interseção sobre
os dez valores é enumerável, o \Cref{thm:borel} segue.

Vale registrar o cálculo explícito da variância, porque ele mostra por que a
aplicação direta de Chebyshev seguida de Borel-Cantelli não conclui.
Para cada $j$,
\[
\operatorname{Var}(Y_j)
=
\mathbb{P}(X_j=k)\bigl(1-\mathbb{P}(X_j=k)\bigr)
=
\frac{1}{10}\cdot\frac{9}{10}
=
\frac{9}{100}.
\]
Como as $Y_j$ são independentes,
\[
\operatorname{Var}\Bigl(\frac{1}{n}\sum_{j=1}^{n}Y_j\Bigr)
=
\frac{9}{100n}.
\]
A cota de Chebyshev é, portanto,
\[
\mathbb{P}\bigl(|\nu_{n,k}-\tfrac{1}{10}|\geqslant\varepsilon\bigr)
\leqslant
\frac{9}{100n\varepsilon^2},
\]
que tende a zero e dá a convergência em probabilidade, mas não é somável. É
precisamente a passagem ao quarto momento que produz a somabilidade, como no
\Cref{thm:quarto-momento}. Esse é o conteúdo do exercício sobre as frequências
dos dígitos proposto nas notas de aula sobre modos de convergência.

O argumento acima demonstrou o \Cref{thm:borel} sem usar a versão geral da lei
forte. A prova vale porque as indicadoras são limitadas. Para uma sequência
geral de variáveis independentes e identicamente distribuídas com
$\mathbb{E}[|X_1|]<\infty$ e sem quarto momento, o argumento não se aplica;
para tratar esse caso, usaremos o método da
\Cref{sec:etemadi}.

\subsection{Dos Dígitos aos Números Reais}\label{subsec:digitos-reais}

Na subseção anterior, demonstramos que uma sequência de dígitos sorteados de
forma independente, com os dez dígitos equiprováveis, tem, quase certamente,
frequência $1/10$ para cada dígito. Queremos agora obter uma afirmação sobre
os dígitos de um número real escolhido ao acaso. Para isso, precisamos
justificar a ligação entre os dois experimentos: sortear os dígitos e
escolher um ponto do intervalo.

A ligação vem da seguinte observação: \emph{fixar os primeiros $n$
dígitos decimais de um número equivale a localizá-lo em um intervalo de
comprimento $10^{-n}$}. Vamos primeiro entender essa observação. Em seguida,
usaremos comprimentos de intervalos para verificar a independência dos
dígitos e aplicar a lei forte.

\subsubsection{Um número escolhido uniformemente}

Denotaremos por $\lambda$ a medida de Lebesgue, que associa a cada intervalo
seu comprimento. Como $\lambda([0,1))=1$, sua restrição a $[0,1)$ é uma medida
de probabilidade. Escolher um ponto uniformemente nesse intervalo significa
justamente usar essa probabilidade: a chance de o ponto escolhido pertencer
a $[a,b)\subseteq[0,1)$ é
\[
\lambda([a,b))=b-a.
\]
Trabalharemos primeiro em $[0,1)$; ao final, incluiremos o extremo $1$.

Para cada $x\in[0,1)$, escrevemos
\[
x=0{,}d_1(x)d_2(x)d_3(x)\ldots
=\sum_{j=1}^{\infty}\frac{d_j(x)}{10^j},
\qquad d_j(x)\in\{0,1,\ldots,9\}.
\]
Aqui, $d_j(x)$ é o dígito que ocupa a $j$-ésima posição depois da vírgula.
Adotamos sempre a expansão que não tem uma cauda infinita de noves. Por
exemplo, entre as duas representações
\[
0{,}37000\ldots=0{,}36999\ldots,
\]
usaremos a primeira. Essa convenção determina os dígitos de cada ponto sem
ambiguidade.

\subsubsection{Dígitos prescritos e intervalos decimais}

Comecemos com o primeiro dígito. Dividindo $[0,1)$ em dez intervalos de mesmo
comprimento, obtemos
\[
\left[0,\frac{1}{10}\right),\quad
\left[\frac{1}{10},\frac{2}{10}\right),\quad\ldots,\quad
\left[\frac{9}{10},1\right).
\]
Esses intervalos correspondem, nessa ordem, aos dez valores possíveis de
$d_1(x)$. Em particular, $d_1(x)=3$ se, e somente se,
$x\in\left[\frac{3}{10},\frac{4}{10}\right)$. Para determinar o segundo dígito, dividimos cada um
desses intervalos em dez partes iguais, e repetimos o procedimento para os
dígitos seguintes. Por exemplo,
\begin{gather*}
\{x:d_1(x)=3\}
=\left[\frac{3}{10},\frac{4}{10}\right),\\[0.4cm]
\{x:d_1(x)=3,\ d_2(x)=7\}
=\left[\frac{37}{100},\frac{38}{100}\right),\\[0.4cm]
\{x:d_1(x)=3,\ d_2(x)=7,\ d_3(x)=2\}
=\left[\frac{372}{1000},\frac{373}{1000}\right).
\end{gather*}
Cada novo dígito especificado reduz o comprimento do intervalo por um fator
$10$, como ilustra a \Cref{fig:intervalos-decimais}. Os extremos esquerdos
estão incluídos e os direitos estão excluídos por causa da convenção adotada:
por exemplo, $0{,}38=0{,}38000\ldots$ tem segundo dígito $8$ e, portanto, não
pertence ao intervalo dos números que começam por $0{,}37$.

\begin{figure}[H]
\centering
\begin{tikzpicture}[x=1.04cm,y=1cm]
  % Em cada linha, o intervalo da linha anterior é ampliado e dividido em dez.
  \foreach \y/\indice in {0/3,-1.8/7,-3.6/2} {
    \fill[border!35] (0,\y) rectangle (10,{\y+0.4});
    \fill[teal!20] (\indice,\y) rectangle ({\indice+1},{\y+0.4});
    \draw[muted,thin] (0,\y) rectangle (10,{\y+0.4});
    \foreach \k in {1,...,9} {
      \draw[muted,thin] (\k,\y) -- (\k,{\y+0.4});
    }
    \foreach \k in {0,...,9} {
      \ifnum\k=\indice\relax
        \node[teal!70!black,font=\small\bfseries]
          at ({\k+0.5},{\y+0.2}) {\k};
      \else
        \node[muted,font=\small] at ({\k+0.5},{\y+0.2}) {\k};
      \fi
    }
    \fill[teal] (\indice,\y) circle (1.7pt);
    \draw[teal,fill=white,line width=0.8pt]
      ({\indice+1},\y) circle (1.7pt);
  }

  % Os rótulos abaixo de cada linha são os valores dos pontos do intervalo.
  \node[below,font=\small] at (0,0) {$0$};
  \node[below,font=\small] at (10,0) {$1$};
  \node[below,font=\small,teal!70!black] at (3,0) {$0{,}3$};
  \node[below,font=\small,teal!70!black] at (4,0) {$0{,}4$};
  \node[right,align=left,font=\small] at (10.65,0.2)
    {$d_1=3$\\comprimento $10^{-1}$};

  \draw[muted!65,dashed] (3,-0.5) -- (0,-1.25);
  \draw[muted!65,dashed] (4,-0.5) -- (10,-1.25);
  \node[below,font=\small] at (0,-1.8) {$0{,}3$};
  \node[below,font=\small] at (10,-1.8) {$0{,}4$};
  \node[below,font=\small,teal!70!black] at (7,-1.8) {$0{,}37$};
  \node[below,font=\small,teal!70!black] at (8,-1.8) {$0{,}38$};
  \node[right,align=left,font=\small] at (10.65,-1.6)
    {$d_1=3,\ d_2=7$\\comprimento $10^{-2}$};

  \draw[muted!65,dashed] (7,-2.3) -- (0,-3.05);
  \draw[muted!65,dashed] (8,-2.3) -- (10,-3.05);
  \node[below,font=\small] at (0,-3.6) {$0{,}37$};
  \node[below,font=\small] at (10,-3.6) {$0{,}38$};
  \node[below,font=\small,teal!70!black] at (2,-3.6) {$0{,}372$};
  \node[below,font=\small,teal!70!black] at (3,-3.6) {$0{,}373$};
  \node[right,align=left,font=\small] at (10.65,-3.4)
    {$d_1=3,\ d_2=7,\ d_3=2$\\comprimento $10^{-3}$};
\end{tikzpicture}
\caption{As linhas usam escalas distintas: o intervalo destacado em cada
linha é ampliado na seguinte. Os números nas dez partes indicam o próximo
dígito; os pontos cheios e vazios indicam se os extremos estão incluídos ou excluídos.}
\label{fig:intervalos-decimais}
\end{figure}

O mesmo raciocínio vale para quaisquer dígitos $u_1,\ldots,u_n$. O número
\[
a=\frac{u_1}{10}+\frac{u_2}{10^2}+\cdots+\frac{u_n}{10^n}
\]
é o extremo esquerdo do intervalo correspondente. Como fizemos $n$ divisões
sucessivas em dez partes iguais, esse intervalo tem comprimento $10^{-n}$.
Portanto,
\begin{equation}\label{eq:intervalo-decimal}
I(u_1,\ldots,u_n)
\equiv\{x\in[0,1):d_1(x)=u_1,\ldots,d_n(x)=u_n\}
=[a,a+10^{-n}).
\end{equation}
Em particular,
\begin{equation}\label{eq:probabilidade-digitos-reais}
\lambda\bigl(I(u_1,\ldots,u_n)\bigr)=10^{-n}
=\mathbb{P}\bigl(C(u_1,\ldots,u_n)\bigr).
\end{equation}
Assim, uma condição sobre os primeiros dígitos tem a mesma probabilidade nos
dois experimentos: no espaço de sequências, ela determina um cilindro; no
intervalo $[0,1)$, ela determina um intervalo decimal.

\subsubsection{Da independência à frequência dos dígitos}

Para aplicar a lei forte às funções $d_j$, precisamos verificar que elas são
variáveis aleatórias independentes e identicamente distribuídas. A
aleatoriedade está na escolha de $x$: depois de escolhido o ponto, todos os
seus dígitos ficam determinados. A igualdade
\eqref{eq:probabilidade-digitos-reais} já trata dos primeiros $n$ dígitos.
A independência exige também considerar posições que não sejam consecutivas.
Podemos fazer isso deixando livres os dígitos das posições não especificadas.

\begin{lemma}[Os dígitos de um ponto uniforme]\label{lem:digitos-uniformes}
No intervalo $[0,1)$ com a probabilidade $\lambda$, as funções $d_1,d_2,\ldots$
são variáveis aleatórias independentes e identicamente distribuídas, com
\[
\lambda\bigl(\{x:d_j(x)=k\}\bigr)=\frac{1}{10}
\qquad\text{para todo }j\geqslant1\text{ e }k\in\{0,1,\ldots,9\}.
\]
\end{lemma}

\begin{proof}
Vejamos primeiro um exemplo com uma posição livre. A condição
$d_1(x)=3$, $d_3(x)=2$ permite que o segundo dígito seja qualquer um dos dez
dígitos. Logo, o evento correspondente é a união disjunta dos dez intervalos
\[
I(3,0,2),\quad I(3,1,2),\quad\ldots,\quad I(3,9,2).
\]
Cada um tem comprimento $10^{-3}$. A probabilidade dessa união é, portanto,
$10\cdot10^{-3}=10^{-2}$.

Em geral, fixamos $r$ posições $1\leqslant n_1<\cdots<n_r$ e os valores dos
dígitos nessas posições, digamos $k_1,\ldots,k_r$. Agora, damos uma atenção especial para a última
posição especificada, $N=n_r$. Entre as primeiras $N$ posições, há $r$
fixadas e $N-r$ livres. Como cada posição livre admite dez escolhas, há
$10^{N-r}$ sequências de $N$ dígitos compatíveis com as condições. Cada uma
determina um intervalo de comprimento $10^{-N}$, e esses intervalos são
disjuntos. Somando seus comprimentos, obtemos
\begin{equation}\label{eq:independencia-digitos-reais}
\lambda\bigl(\{x:d_{n_1}(x)=k_1,\ldots,d_{n_r}(x)=k_r\}\bigr)
=10^{N-r}\cdot10^{-N}=10^{-r}.
\end{equation}
Essa descrição mostra também que os eventos considerados são mensuráveis,
pois são uniões finitas de intervalos. Em particular, cada $d_j$ é uma
variável aleatória, já que assume apenas dez valores e cada conjunto
$\{x:d_j(x)=k\}$ é mensurável.

Tomando $r=1$ em \eqref{eq:independencia-digitos-reais}, vemos que cada
$d_j$ tem distribuição uniforme em $\{0,1,\ldots,9\}$. Para $r$ qualquer, a
mesma igualdade fornece
\[
\lambda\bigl(\{d_{n_1}=k_1,\ldots,d_{n_r}=k_r\}\bigr)
=\prod_{i=1}^{r}\lambda\bigl(\{d_{n_i}=k_i\}\bigr).
\]
Como cada $d_j$ assume um número finito de valores, a fatoração para conjuntos
de valores se obtém somando essas igualdades. Isso prova a independência.
\end{proof}

O lema justifica o uso do modelo de sorteios independentes para estudar os
dígitos de um ponto uniforme. Agora podemos repetir a aplicação da lei forte
feita na subseção anterior, usando a medida $\lambda$ como probabilidade.
Nesse espaço, dizer que uma propriedade vale \emph{quase certamente} significa
que o conjunto dos pontos onde ela falha tem medida de Lebesgue zero. É
exatamente esse o sentido de dizer que ela vale para \emph{quase todo} ponto.

\begin{theorem}[Borel]\label{thm:borel-reais}
Para quase todo $x\in[0,1]$, a frequência relativa de cada dígito na expansão
decimal de $x$ tende a $1/10$. Em outras palavras, quase todo número real é
simplesmente normal na base $10$.
\end{theorem}

\begin{proof}
Fixe um dígito $k\in\{0,1,\ldots,9\}$. Para contar suas ocorrências,
definimos, no espaço $[0,1)$,
\[
Z_j(x)=\mathds{1}_{\{d_j=k\}}(x)
=\begin{cases}
1,&\text{se }d_j(x)=k,\\
0,&\text{se }d_j(x)\neq k.
\end{cases}
\]
Pelo \Cref{lem:digitos-uniformes}, as variáveis $Z_j$ são independentes e
identicamente distribuídas. Além disso, são limitadas por $1$ e têm esperança
\[
\int_{[0,1)}Z_j(x)\,d\lambda(x)
=\lambda\bigl(\{x:d_j(x)=k\}\bigr)=\frac{1}{10}.
\]
O \Cref{cor:limitadas} se aplica e fornece
\[
\frac{\#\{1\leqslant j\leqslant n:d_j(x)=k\}}{n}
=\frac{1}{n}\sum_{j=1}^{n}Z_j(x)
\longrightarrow\frac{1}{10}
\qquad\text{para $\lambda$-quase todo }x\in[0,1).
\]
O numerador conta quantas vezes o dígito $k$ aparece nas primeiras $n$
posições; portanto, a razão é precisamente sua frequência relativa.

Para cada $k$, o conjunto excepcional tem medida zero. A união desses dez
conjuntos também tem medida zero, de modo que, fora dela, as dez frequências
convergem simultaneamente para $1/10$. Incluir o ponto $1$ não altera a
afirmação, pois $\lambda(\{1\})=0$. A escolha entre as duas expansões dos
decimais finitos tampouco afeta o resultado: os pontos que admitem duas
expansões formam um conjunto enumerável e, por isso, de medida zero.

Finalmente, para um número real $x$, consideramos os dígitos de sua parte
fracionária $x-\lfloor x\rfloor\in[0,1)$. Se $N\subseteq[0,1)$ é o conjunto
excepcional obtido acima, o conjunto excepcional em toda a reta é
\[
\bigcup_{m\in\mathbb{Z}}(m+N).
\]
Cada translação $m+N$ tem medida zero, e a união é enumerável. Isso prova a
afirmação em $\mathbb{R}$.
\end{proof}

%----------------------------------------------------------------------
\section{A Prova de Etemadi da Lei Forte}\label{sec:etemadi}

\subsection{Enunciado e Visão Geral}\label{subsec:etemadi-enunciado}

O resultado central destas notas é o seguinte teorema, devido a Etemadi. Ele
estabelece a lei forte na hipótese mínima sobre os momentos e necessita apenas
da hipótese mais fraca de independência dois a dois, e não de independência
mútua.

\begin{theorem}[Lei forte dos grandes números]\label{thm:lei-forte}
Sejam $X_1,X_2,\ldots$ variáveis aleatórias independentes dois a dois e
identicamente distribuídas, com $\mathbb{E}[|X_i|]<\infty$. Seja
$\mu=\mathbb{E}[X_i]$ e $S_n=X_1+\cdots+X_n$. Então
\[
\frac{S_n}{n}\longrightarrow\mu
\qquad\text{quase certamente.}
\]
\end{theorem}

Antes de apresentar os detalhes da prova, vale entender como ela amplia o
resultado anterior. O \Cref{thm:quarto-momento} exige
$\mathbb{E}[X_1^4]<\infty$, enquanto o \Cref{thm:lei-forte} requer apenas
$\mathbb{E}[|X_1|]<\infty$. Para ver concretamente a diferença entre essas
hipóteses, considere uma variável aleatória $X_1$ com densidade
\[
f(x)=\begin{cases}
\dfrac{3}{2x^{5/2}},&\text{se }x\geqslant1,\\[0.5cm]
0,&\text{se }x<1.
\end{cases}
\]
Essa é a distribuição de Pareto com escala $1$ e parâmetro de forma $3/2$,
que estudaremos na \Cref{subsec:pareto}. Abaixo mostramos o esboço do gráfico desta função densidade na
\Cref{fig:densidade-pareto-exemplo}.

\begin{figure}[H]
\centering
\begin{tikzpicture}[x=1.45cm,y=2.8cm]
  \draw[thin,-stealth] (-0.7,0) -- (6.5,0) node[below right] {$x$};
  \draw[thin,-stealth] (0,-0.08) -- (0,1.8) node[left] {$f(x)$};
  \node[below left,font=\small] at (0,0) {$0$};
  \foreach \x in {1,...,6} {
    \draw[thin] (\x,0.02) -- (\x,-0.02)
      node[below=2pt,font=\small] {$\x$};
  }
  \draw[gray!65,thin,dashed] (0,1.5) -- (1,1.5);
  \draw[gray!65,thin,dashed] (1,0) -- (1,1.5);
  \node[left,font=\small] at (0,1.5) {$\frac{3}{2}$};

  % Os dois ramos são desenhados separadamente para representar o salto em 1.
  \draw[teal,thick] (-0.7,0) -- (1,0);
  \draw[teal,thick,domain=1:6.2,samples=180,smooth]
    plot (\x,{3/(2*\x^2.5)});
  \draw[teal,fill=white,line width=0.8pt] (1,0) circle (1.7pt);
  \fill[teal] (1,1.5) circle (1.7pt);
  \node[font=\small,anchor=west] at (2.2,0.95)
    {$f(x)=\dfrac{3}{2x^{5/2}},\quad x\geqslant1$};
\end{tikzpicture}
\caption{Densidade de Pareto com escala $1$ e parâmetro de forma $3/2$.}
\label{fig:densidade-pareto-exemplo}
\end{figure}

Como $X_1\geqslant1$ quase certamente,
o cálculo dos dois momentos fornece
\begin{gather*}
\mathbb{E}[|X_1|]
=\int_1^{\infty}x\,\frac{3}{2x^{5/2}}\,dx
=\frac{3}{2}\int_1^{\infty}x^{-3/2}\,dx=3<\infty,\\[0.4cm]
\mathbb{E}[X_1^4]
=\int_1^{\infty}x^4\,\frac{3}{2x^{5/2}}\,dx
=\frac{3}{2}\int_1^{\infty}x^{3/2}\,dx=\infty.
\end{gather*}
Portanto, para uma sequência independente e identicamente distribuída com
essa distribuição, o \Cref{thm:lei-forte} garante a convergência quase certa
das médias amostrais para $3$, embora a hipótese do
\Cref{thm:quarto-momento} não seja satisfeita.


\bigskip

A prova de Etemadi substitui o quarto momento por uma
estimativa de momentos das variáveis truncadas, cujo nível de truncamento
cresce com o índice. Com apenas a integrabilidade, não temos garantia de que a
cota de Chebyshev seja somável em todos os índices. A estimativa abaixo
garante a somabilidade ao longo de índices de crescimento geométrico, e a
monotonicidade permite controlar os índices restantes.

A prova tem quatro etapas, que enunciamos antes de executá-las:

\begin{enumerate}
\item[(i)] \emph{Redução a variáveis não negativas.} Decompomos
$X_n=X_n^+-X_n^-$ e tratamos cada parte separadamente.
\item[(ii)] \emph{Truncamento.} Substituímos $X_k$ por
$Y_k=X_k\mathds{1}_{\{|X_k|\leqslant k\}}$ e mostramos que basta provar a lei
forte para as variáveis truncadas; em seguida, identificamos o limite das
médias de suas esperanças.
\item[(iii)] \emph{Estimativa de momentos.} Mostramos que
$\sum_k\operatorname{Var}(Y_k)/k^2<\infty$.
\item[(iv)] \emph{Subsequência geométrica e monotonicidade.} Aplicamos
Chebyshev ao longo dos índices $k(n)=\lfloor\alpha^n\rfloor$, concluímos por
Borel-Cantelli, e usamos a monotonicidade da soma truncada para controlar os
índices intermediários.
\end{enumerate}

No livro de Durrett, a redução a variáveis não negativas aparece depois das
estimativas de momentos. Nestas notas, apresentamos essa redução primeiro
para destacar a propriedade que será usada na etapa final: as somas parciais
de variáveis não negativas são não decrescentes. Essa monotonicidade permitirá
estender a convergência das médias, obtida na subsequência geométrica,
a todos os índices.

\subsection{Redução a Variáveis Não Negativas}\label{subsec:reducao}

Para $x\in\mathbb{R}$, escrevemos $x^+=\max\{x,0\}$ e $x^-=\max\{-x,0\}$, de
modo que $x=x^+-x^-$ e $|x|=x^++x^-$. Definimos
\[
X_n^+=\max\{X_n,0\},
\qquad
X_n^-=\max\{-X_n,0\}.
\]
As duas famílias $(X_n^+)$ e $(X_n^-)$ são formadas por variáveis não
negativas, independentes dois a dois e identicamente distribuídas. Além disso,
$0\leqslant X_n^\pm\leqslant|X_n|$, logo
$\mathbb{E}[X_n^\pm]\leqslant\mathbb{E}[|X_n|]<\infty$.

Suponha que o \Cref{thm:lei-forte} esteja demonstrado para variáveis não
negativas. Aplicando-o a cada família, obtemos, quase certamente,
\[
\frac{1}{n}\sum_{k=1}^{n}X_k^+\longrightarrow\mathbb{E}[X_1^+],
\qquad
\frac{1}{n}\sum_{k=1}^{n}X_k^-\longrightarrow\mathbb{E}[X_1^-].
\]
Os dois conjuntos de convergência têm probabilidade $1$, logo a interseção
também tem probabilidade $1$. Em cada ponto dessa interseção, subtraindo as
os dois limites, obtemos
\[
\frac{S_n}{n}
=
\frac{1}{n}\sum_{k=1}^{n}X_k^+
-
\frac{1}{n}\sum_{k=1}^{n}X_k^-
\longrightarrow
\mathbb{E}[X_1^+]-\mathbb{E}[X_1^-]
=
\mathbb{E}[X_1]
=
\mu.
\]
Portanto, basta provar o teorema sob a hipótese adicional $X_n\geqslant0$ para
todo $n$. Nesse caso, $\mu=\mathbb{E}[X_1]\geqslant0$ e as somas truncadas
serão não decrescentes, o que será usado na \Cref{subsec:intermediarios}.

\subsection{O Truncamento}\label{subsec:lema-truncamento}

Pelo argumento dado acima, podemos supor que $X_k\geqslant0$ para
todo $k$. É claro que a hipótese de integrabilidade não é suficiente para garantir que essas variáveis
tenham segundo momento finito, uma condição necessária para obter uma cota
finita pela desigualdade de Chebyshev. Por isto vamos substituir cada v.a.
$X_k$ por uma variável aleatória limitada que será denotada por $Y_k$
que é definida a partir de um truncamento de $X_k$ e as somas
parciais destas novas variáveis serão denotadas por $T_n$. Mais precisamente, definimos
\begin{equation}\label{eq:truncamento}
Y_k\equiv X_k\cdot \mathds{1}_{\{|X_k|\leqslant k\}},
\qquad
T_n\equiv Y_1+\cdots+Y_n.
\end{equation}
Assim, $Y_k$ coincide com $X_k$ ao longo de $\{X_k\leqslant k\}$ e vale zero
no evento $\{X_k>k\}$. Como $0\leqslant Y_k\leqslant k$, temos
$\mathbb{E}[Y_k^2]\leqslant k^2<\infty$. A independência dois a dois também
é preservada, pois cada $Y_k$ é uma função boreliana da respectiva variável
$X_k$. Como o nível de truncamento varia com $k$, as variáveis $Y_k$ não
são, em geral, identicamente distribuídas.

A escolha do nível $k$ é motivado pelo critério de comparação
\eqref{eq:comparacao-cauda} que está relacionado a integrabilidade de $X_1$ via a convergência
da série $\sum_{k=1}^{\infty}\mathbb{P}(X_1>k)$. Como $X_k$ e $X_1$ têm a
mesma distribuição, temos
\[
\mathbb{P}(X_k\neq Y_k)
=\mathbb{P}(X_k>k)
=\mathbb{P}(X_1>k).
\]
O primeiro lema de Borel-Cantelli
permitirá mostrar que, quase certamente, o truncamento modifica apenas um
número finito de termos. A contribuição dessas modificações para a
diferença entre as médias tende a zero, como afirma o lema.

\begin{lemma}[Truncamento]\label{lem:truncamento}
Assuma que são válidas as hipóteses do \Cref{thm:lei-forte} e sejam $Y_k$ e $T_n$ as variáveis aleatórias
definidos em \eqref{eq:truncamento}. Então
\[
\frac{S_n-T_n}{n}\longrightarrow0
\qquad\text{quase certamente.}
\]
Em particular, $S_n/n\to\mu$ quase certamente se, e somente se,
$T_n/n\to\mu$ quase certamente.
\end{lemma}

\begin{proof}
Vamos mostrar que, quase certamente, o truncamento altera apenas um número
finito de termos. Os eventos correspondentes são
\[
A_k\equiv\{X_k\neq Y_k\}=\{|X_k|>k\}.
\]
Como as variáveis são identicamente distribuídas,
$\mathbb{P}(|X_k|>k)=\mathbb{P}(|X_1|>k)$. A fórmula de esperança por integração
da cauda e a comparação \eqref{eq:comparacao-cauda} dão a seguinte estimativa
(a soma é de termos não negativos):
\begin{equation}\label{eq:serie-truncamento}
\sum_{k=1}^{\infty}\mathbb{P}(|X_k|>k)
=
\sum_{k=1}^{\infty}\mathbb{P}(|X_1|>k)
\leqslant
\int_0^{\infty}\mathbb{P}(|X_1|>t)\,dt
=
\mathbb{E}[|X_1|].
\end{equation}
Por hipótese, o lado direito é finito.

Por \eqref{eq:serie-truncamento} e pelo primeiro lema de Borel-Cantelli,
\[
\mathbb{P}\bigl(\limsup_{k\to\infty}A_k\bigr)=0.
\]
Fora desse conjunto nulo, o conjunto dos índices $k$ para os quais
$X_k\neq Y_k$ é finito. Para tais $\omega$, definimos
\[
K(\omega)=\max\{k:X_k(\omega)\neq Y_k(\omega)\},
\qquad
R(\omega)=\sum_{k=1}^{K(\omega)}\bigl|X_k(\omega)-Y_k(\omega)\bigr|,
\]
com $K(\omega)=0$ e $R(\omega)=0$ quando o conjunto é vazio. Como cada
diferença é finita e o máximo é tomado sobre um conjunto finito,
$R(\omega)<\infty$.

Para $k>K(\omega)$ temos $Y_k(\omega)=X_k(\omega)$. Logo, para todo
$n\geqslant1$,
\[
\bigl|S_n(\omega)-T_n(\omega)\bigr|
=
\Bigl|\sum_{k=1}^{n}\bigl(X_k(\omega)-Y_k(\omega)\bigr)\Bigr|
\leqslant
\sum_{k=1}^{K(\omega)}\bigl|X_k(\omega)-Y_k(\omega)\bigr|
=
R(\omega).
\]
Dividindo por $n$,
\[
\Bigl|\frac{S_n(\omega)}{n}-\frac{T_n(\omega)}{n}\Bigr|
\leqslant
\frac{R(\omega)}{n}\longrightarrow0.
\]
Isso prova que $(S_n-T_n)/n\to0$ quase certamente. Consequentemente,
$S_n/n\to\mu$ quase certamente se, e somente se, $T_n/n\to\mu$ quase
certamente.
\end{proof}

A cota obtida é uniforme em $n$, e é isso que permite dividir por $n$ sem
discutir subsequências. O ponto é que $K(\omega)$ é o \emph{último} índice em
que o truncamento interfere; depois dele, as diferenças são exatamente nulas e
a soma parcial para de crescer. Uma cota dependente de $n$ exigiria verificar
adicionalmente que ela tende a zero após a divisão por $n$.

\subsection{As Médias Truncadas}\label{subsec:medias-truncadas}

Antes de estimar as flutuações de $T_n$, identificamos o limite de suas médias.
Como já reduzimos ao caso $X_k\geqslant0$, a variável
\[
Z_k=X_1\mathds{1}_{\{X_1\leqslant k\}}
\]
cresce para $X_1$ quando $k\to\infty$. Pelo Teorema da Convergência Monótona,
\[
\mathbb{E}[Z_k]\longrightarrow\mathbb{E}[X_1]=\mu.
\]
Para transferir essa convergência para as esperanças de $Y_k$, vamos mostrar
que, para cada $k$ fixado, as variáveis
\[
Y_k=X_k\mathds{1}_{\{X_k\leqslant k\}}
\qquad\text{e}\qquad
Z_k=X_1\mathds{1}_{\{X_1\leqslant k\}}
\]
têm a mesma distribuição. Nas duas definições, usamos o mesmo nível de
truncamento $k$. Assim, aplicamos a mesma regra a $X_k$ e a $X_1$, que têm a
mesma distribuição. Para verificar que $Y_k$ e $Z_k$ também têm a mesma
distribuição, basta comparar suas funções de distribuição.

Como $0\leqslant Y_k\leqslant k$ e $0\leqslant Z_k\leqslant k$, para $t<0$
temos $\mathbb{P}(Y_k\leqslant t)=\mathbb{P}(Z_k\leqslant t)=0$, e para
$t\geqslant k$ essas duas probabilidades são iguais a $1$.
Resta considerar $0\leqslant t<k$. Nesse caso, $Y_k\leqslant t$ ocorre
quando $X_k\leqslant t$ ou quando $X_k>k$, pois, nesta última situação,
o truncamento atribui a $Y_k$ o valor zero. Portanto,
\[
\{Y_k\leqslant t\}
=\{X_k\leqslant t\}\cup\{X_k>k\}.
\]
Os dois eventos à direita são disjuntos, e a mesma decomposição vale para
$Z_k$, substituindo $X_k$ por $X_1$. Como $X_k$ e $X_1$ têm a mesma
distribuição, obtemos
\begin{align*}
\mathbb{P}(Y_k\leqslant t)
&=\mathbb{P}(X_k\leqslant t)+\mathbb{P}(X_k>k)
\\[0.3cm]
&=\mathbb{P}(X_1\leqslant t)+\mathbb{P}(X_1>k)
\\[0.3cm]
&=\mathbb{P}(Z_k\leqslant t).
\end{align*}
Assim, as funções de distribuição de $Y_k$ e $Z_k$ coincidem em todo
$t\in\mathbb{R}$, o que prova a igualdade em distribuição. Como essas
variáveis são limitadas, suas esperanças são finitas e iguais. Logo,
\[
\mathbb{E}[Y_k]=\mathbb{E}[Z_k]\longrightarrow\mu.
\]
Usaremos agora a seguinte propriedade das sequências numéricas: se
$a_m\to a$, então suas médias aritméticas também convergem para $a$, isto é,
\[
\frac{a_1+\cdots+a_n}{n}\longrightarrow a.
\]
Essa propriedade é conhecida como convergência de Cesàro. A ideia é que a
contribuição de um número finito de termos iniciais para a média tende a zero
após a divisão por $n$; a partir de certo índice, todos os demais termos
estão tão próximos de $a$ quanto quisermos.

Aplicamos essa propriedade à sequência $a_m=\mathbb{E}[Y_m]$, cujo limite
é $\mu$. Pela linearidade da esperança, obtemos
\begin{equation}\label{eq:cesaro}
\frac{\mathbb{E}[T_n]}{n}
=
\frac{1}{n}\sum_{m=1}^{n}\mathbb{E}[Y_m]
\longrightarrow\mu.
\end{equation}
Essa identificação vale para a sequência inteira; mais adiante, será usada
também nos índices da subsequência geométrica.

\subsection{A Estimativa de Momentos}\label{subsec:estimativa-momentos}

Na subseção anterior, mostramos que $\mathbb{E}[T_n]/n\to\mu$.
Resta controlar a diferença entre a média truncada $T_n/n$ e sua esperança,
isto é, provar que $(T_n-\mathbb{E}[T_n])/n\to0$ quase certamente.
Como as variáveis aleatórias $Y_k$ têm segundo momento finito e são duas-a-duas independentes,
a variância de $T_n$ é a soma das variâncias dos seus termos. Portanto,
para todo $\varepsilon>0$, a desigualdade de Chebyshev fornece
\begin{align}\label{eq-aux1}
\mathbb{P}\left(
\left|\frac{T_n-\mathbb{E}[T_n]}{n}\right|>\varepsilon
\right)
\leqslant
\frac{1}{\varepsilon^2n^2}
\sum_{k=1}^{n}\operatorname{Var}(Y_k).
\end{align}
A expressão acima sugere que a prova do teorema poderia
ser finalizada se fosse possível aplicar o critério de Borel-Cantelli.
O problema é justamente encontrar uma cota superior somável para a expressão que
aparece no lado direito da desigualdade \eqref{eq-aux1}.
A grande ideia de Etemadi foi perceber dois fatos fundamentais: que a convergência quase certa da sequência
\begin{align}\label{eq-aux2}
\frac{T_{k(n)}}{k(n)},
\end{align}
onde $k(n)$ é uma sequência apropriada de índices, implica na convergência quase certa
de toda a sequência $T_n/n$ e; que a convergência quase certa da sequência \eqref{eq-aux2} poderia
ser obtida por \eqref{eq-aux1}, escolhendo $k(n)$ apropriadamente.

Antes de desenvolver a ideia de Etemadi precisamos de alguns lemas técnicos.

\begin{lemma}\label{lem:harmonica}
Para todo $y\geqslant0$,
\[
2y\sum_{k>y}\frac{1}{k^2}\leqslant4.
\]
\end{lemma}

\begin{proof}
Começamos pela comparação com a integral. Se $m\geqslant2$, a função
$x\mapsto x^{-2}$ é decrescente em $[m-1,\infty)$. Para $x\in[k-1,k]$ com
$k\geqslant m$, temos $x\leqslant k$, logo $x^{-2}\geqslant k^{-2}$ e
\[
\frac{1}{k^2}\leqslant\int_{k-1}^{k}\frac{1}{x^2}\,dx.
\]
Somando em $k\geqslant m$,
\begin{equation}\label{eq:cota-harmonica}
\sum_{k\geqslant m}\frac{1}{k^2}
\leqslant
\int_{m-1}^{\infty}\frac{dx}{x^2}
=
\frac{1}{m-1}.
\end{equation}

O restante da prova é divido em dois casos. Suponha inicialmente que $y\geqslant1$.
Neste caso, a soma $\sum_{k>y}k^{-2}$ começa em
$k=\lfloor y\rfloor+1\geqslant2$, de modo que \eqref{eq:cota-harmonica} se
aplica com $m=\lfloor y\rfloor+1$, fornecendo a seguinte desigualdade
\[
2y\sum_{k>y}\frac{1}{k^2}
=
2y\sum_{k=\lfloor y\rfloor+1}^{\infty}\frac{1}{k^2}
\leqslant
\frac{2y}{\lfloor y\rfloor}
\leqslant
4,
\]
pois, para $y\geqslant1$, temos $y<\lfloor y\rfloor+1\leqslant2\lfloor
y\rfloor$.

Para finalizar resta obter a desigualdade para o caso $0\leqslant y<1$.
Neste caso, temos que todo inteiro positivo $k$ satisfaz
$k>y$, e portanto segue de uma aplicação diretta de \eqref{eq:cota-harmonica} com $m=2$ que
\[
2y\sum_{k>y}\frac{1}{k^2}
=
2y\sum_{k\geqslant1}\frac{1}{k^2}
\leqslant
2\sum_{k\geqslant1}\frac{1}{k^2}
=
2\Bigl(1+\sum_{k\geqslant2}\frac{1}{k^2}\Bigr)
\leqslant
2(1+1)
=
4.
\]
Em ambos os casos a constante $4$ serve. A cota é grosseira, mas suficiente.
\end{proof}

Observe que a constante exata não importa: qualquer constante finita serviria,
porque ela será multiplicada por $\mathbb{E}[|X_1|]$, que é finito. O que importa
é a ordem de grandeza, isto é, o fato de a constante não depender de $y$.

\begin{lemma}[Estimativa de momentos]\label{lem:momento}
Se $Y_k$, para cada $k\in\mathbb{N}$, é a variável aleatória definida em \eqref{eq:truncamento},
então
\begin{align}\label{eq-aux3}
\sum_{k=1}^{\infty}\frac{\operatorname{Var}(Y_k)}{k^2}
\leqslant
4\,\mathbb{E}[|X_1|]
<
\infty.
\end{align}
\end{lemma}

\begin{proof}
Como $\operatorname{Var}(Y_k)=\mathbb{E}[Y_k^2]-(\mathbb{E}[Y_k])^2\leqslant
\mathbb{E}[Y_k^2]$, basta cotar a soma dos segundos momentos e
para estimar $\mathbb{E}[Y_k^2]$, vamos usar a representação do segundo momento mostrada abaixo. Para
qualquer variável $W$ com $\mathbb{E}[W^2]<\infty$,
\begin{equation}\label{eq:camadas}
\mathbb{E}[W^2]
=
\int_0^{\infty}2y\,\mathbb{P}(|W|>y)\,dy.
\end{equation}

Como $|Y_k|\leqslant k$, temos
$\mathbb{P}(|Y_k|>y)=0$ para $y\geqslant k$. Para $0\leqslant y<k$, o evento
$\{|Y_k|>y\}$ está contido em $\{|X_k|>y\}$, pois $|Y_k|\leqslant|X_k|$. Logo
\[
\mathbb{P}(|Y_k|>y)\leqslant\mathds{1}_{\{y<k\}}\,\mathbb{P}(|X_k|>y)
=
\mathds{1}_{\{y<k\}}\,\mathbb{P}(|X_1|>y),
\]
onde na última igualdade usamos que $(X_n)_{n\in\mathbb{N}}$ é uma sequência de v.a.'s identicamente distribuídas.
Portanto
\begin{equation}\label{eq:cota-Yk}
\mathbb{E}[Y_k^2]
\leqslant
\int_0^{k}2y\,\mathbb{P}(|X_1|>y)\,dy.
\end{equation}

Substituindo \eqref{eq:cota-Yk} na soma \eqref{eq-aux3} e trocando a ordem de soma e integral,
o que é permitido pelo teorema de Tonelli porque todos os termos são não
negativos, obtemos a seguinte estimativa
\begin{align*}
\sum_{k=1}^{\infty}\frac{\mathbb{E}[Y_k^2]}{k^2}
&\leqslant
\sum_{k=1}^{\infty}\frac{1}{k^2}\int_0^{\infty}
\mathds{1}_{\{y<k\}}\,2y\,\mathbb{P}(|X_1|>y)\,dy
\\[0.3cm]
&=
\int_0^{\infty}
\Bigl\{\sum_{k=1}^{\infty}\frac{\mathds{1}_{\{y<k\}}}{k^2}\Bigr\}
\,2y\,\mathbb{P}(|X_1|>y)\,dy
\\[0.3cm]
&=
\int_0^{\infty}
\Bigl\{\sum_{k>y}\frac{1}{k^2}\Bigr\}
\,2y\,\mathbb{P}(|X_1|>y)\,dy.
\end{align*}
A soma entre chaves é cotada pelo \Cref{lem:harmonica}
\[
2y\sum_{k>y}\frac{1}{k^2}\leqslant4
\qquad\text{para todo }y\geqslant0.
\]
Logo concluímos das desigualdades acima que
\[
\sum_{k=1}^{\infty}\frac{\operatorname{Var}(Y_k)}{k^2}
\leqslant
\sum_{k=1}^{\infty}\frac{\mathbb{E}[Y_k^2]}{k^2}
\leqslant
\int_0^{\infty}4\,\mathbb{P}(|X_1|>y)\,dy
=
4\,\mathbb{E}[|X_1|]
<
\infty,
\]
onde na última igualdade usamos novamente a fórmula da esperança por integração da cauda
e a hipótese de finitude da esperança de $|X_1|$ e com isto encerramos a prova do lema.
\end{proof}

\subsection{O Argumento de Subsequência}\label{subsec:subsequencia}

Suponha, como permite a \Cref{subsec:reducao}, que $X_n\geqslant0$ para todo
$n$. Então $Y_n\geqslant0$ e $\mu=\mathbb{E}[X_1]\geqslant0$.

Seja $\alpha>1$ fixo, porém arbitrário. Defina
\begin{equation}\label{eq:kn}
k(n)=\lfloor\alpha^n\rfloor,
\qquad n\geqslant1.
\end{equation}
A sequência $k(n)$ é não decrescente e tende a infinito, podendo repetir
alguns termos iniciais.

Queremos aplicar o primeiro lema de Borel-Cantelli às médias truncadas nesses
índices. Para isso, precisamos mostrar que as probabilidades de desvio
obtidas pela desigualdade de Chebyshev têm soma finita. Mais precisamente,
provaremos que, para cada $\varepsilon>0$,
\[
\sum_{n=1}^{\infty}
\mathbb{P}\bigl(|T_{k(n)}-\mathbb{E}[T_{k(n)}]|\geqslant\varepsilon k(n)\bigr)
\leqslant
C_{\alpha,\varepsilon}
\sum_{m=1}^{\infty}\frac{\operatorname{Var}(Y_m)}{m^2},
\]
onde $C_{\alpha,\varepsilon}$ é uma constante finita que depende apenas de
$\alpha$ e de $\varepsilon$. O \Cref{lem:momento} garante que a série do
lado direito converge. Essa estimativa fornecerá, portanto, a somabilidade
necessária para aplicar Borel-Cantelli.

A escolha dos índices geométricos permite obter essa estimativa porque
$k(n)\geqslant\alpha^n/2$ e, consequentemente,
$k(n)^{-2}\leqslant4\alpha^{-2n}$. Os termos $\alpha^{-2n}$ formam uma
progressão geométrica de razão $\alpha^{-2}<1$, cujas somas podemos calcular
explicitamente. No argumento abaixo, usaremos essa comparação depois de
trocar a ordem das somas, para obter a série
$\sum_{m=1}^{\infty}\operatorname{Var}(Y_m)/m^2$ que aparece no lado direito.


\bigskip


Como os $Y_m$ são v.a.'s duas-a-duas independentes, então
\[
\operatorname{Var}(T_{k(n)})
=
\sum_{m=1}^{k(n)}\operatorname{Var}(Y_m).
\]
Pela Desigualdade de Chebyshev sabemos que para cada $\varepsilon>0$ fixado temos
\[
\mathbb{P}\bigl(|T_{k(n)}-\mathbb{E}[T_{k(n)}]|\geqslant\varepsilon k(n)\bigr)
\leqslant
\frac{\operatorname{Var}(T_{k(n)})}{\varepsilon^2k(n)^2}
=
\frac{1}{\varepsilon^2 k(n)^2}\sum_{m=1}^{k(n)} \operatorname{Var}(Y_m).
\]
Próximo passo é somar ambos os lados da desigualdade acima sobre $n\geqslant1$. O lado direito
passa a ser uma soma dupla, onde para cada $n\in\mathbb{N}$, somamos os termos com
$1\leqslant m\leqslant k(n)$.

Para inverter a ordem das somas, fixemos agora $m\geqslant1$. A parcela
$\operatorname{Var}(Y_m)/k(n)^2$ aparece precisamente para os índices $n$
que satisfazem $k(n)\geqslant m$. Assim, podemos somar primeiro sobre esses
índices e depois sobre todos os $m\geqslant1$. Como todos os termos são não
negativos, o teorema de Tonelli permite essa troca mesmo antes de sabermos
se a soma é finita. Obtemos
\begin{align}
\sum_{n=1}^{\infty}
\mathbb{P}\bigl(|T_{k(n)}-\mathbb{E}[T_{k(n)}]|\geqslant\varepsilon k(n)\bigr)
&\leqslant
\frac{1}{\varepsilon^2}
\sum_{n=1}^{\infty}\sum_{m=1}^{k(n)}
\frac{\operatorname{Var}(Y_m)}{k(n)^2}
\notag\\[0.4cm]
&=
\frac{1}{\varepsilon^2}
\sum_{m=1}^{\infty}
\sum_{\substack{n\geqslant1\\k(n)\geqslant m}}
\frac{\operatorname{Var}(Y_m)}{k(n)^2}
\notag\\[0.4cm]
&=
\frac{1}{\varepsilon^2}
\sum_{m=1}^{\infty}\operatorname{Var}(Y_m)
\sum_{\substack{n\geqslant1\\k(n)\geqslant m}}\frac{1}{k(n)^2}.
\label{eq:troca-somas}
\end{align}
Na última igualdade, colocamos $\operatorname{Var}(Y_m)$ em evidência na
soma em $n$, pois essa variância não depende de $n$.

Próximo passo é dominar a soma interna pela cauda de uma série geométrica. Como
$k(n)=\lfloor\alpha^n\rfloor\geqslant\alpha^n/2$ para $n\geqslant1$, temos
$k(n)^{-2}\leqslant4\alpha^{-2n}$. Defina
\[
n_0(m)=\min\{n\geqslant1:\alpha^n\geqslant m\}.
\]
Como $m$ é inteiro, $\lfloor\alpha^n\rfloor\geqslant m$ equivale a
$\alpha^n\geqslant m$, logo os índices da soma são exatamente
$n\geqslant n_0(m)$. Escrevendo $n_0=n_0(m)$, temos
$\alpha^{-2n_0}\leqslant m^{-2}$ e
\begin{equation}\label{eq:cauda-geometrica}
\sum_{n:\,k(n)\geqslant m}\frac{1}{k(n)^2}
\leqslant
4\sum_{n\geqslant n_0}\alpha^{-2n}
=
4\,\alpha^{-2n_0}\sum_{j=0}^{\infty}\alpha^{-2j}
\leqslant
\frac{4}{1-\alpha^{-2}}\cdot\frac{1}{m^2}.
\end{equation}
Substituindo \eqref{eq:cauda-geometrica} em \eqref{eq:troca-somas} e usando o
\Cref{lem:momento},
\begin{equation}\label{eq:serie-Chebyshev}
\sum_{n=1}^{\infty}
\mathbb{P}\bigl(|T_{k(n)}-\mathbb{E}[T_{k(n)}]|\geqslant\varepsilon k(n)\bigr)
\leqslant
\frac{4}{\varepsilon^2(1-\alpha^{-2})}
\sum_{m=1}^{\infty}\frac{\operatorname{Var}(Y_m)}{m^2}
\leqslant
\frac{16\,\mathbb{E}[|X_1|]}{\varepsilon^2(1-\alpha^{-2})}
<
\infty.
\end{equation}
Como $k(n)\geqslant1$, a estimativa \eqref{eq:serie-Chebyshev} mostra que,
para todo $\varepsilon>0$,
\[
\sum_{n=1}^{\infty}
\mathbb{P}\left(
\left|\frac{T_{k(n)}-\mathbb{E}[T_{k(n)}]}{k(n)}\right|
\geqslant\varepsilon
\right)<\infty.
\]
Está, portanto, verificada a hipótese do critério de Borel-Cantelli para a
convergência quase certa, recordado na \Cref{subsec:bc-criterio}. Aplicando
esse critério à sequência dos desvios normalizados, obtemos
\begin{equation}\label{eq:subsequencia}
\frac{T_{k(n)}-\mathbb{E}[T_{k(n)}]}{k(n)}\longrightarrow0
\qquad\text{quase certamente.}
\end{equation}

Além disso, em \eqref{eq:cesaro} já mostramos que a sequência numérica
$\mathbb{E}[T_m]/m$ converge para $\mu$. Como $k(n)\to\infty$, essa
convergência também vale ao longo dos índices $k(n)$:
\[
\frac{\mathbb{E}[T_{k(n)}]}{k(n)}\longrightarrow\mu.
\]
Para concluir, decompomos a média truncada em duas parcelas:
\[
\frac{T_{k(n)}}{k(n)}
=\frac{T_{k(n)}-\mathbb{E}[T_{k(n)}]}{k(n)}
+\frac{\mathbb{E}[T_{k(n)}]}{k(n)}.
\]
A primeira tende a zero quase certamente, por \eqref{eq:subsequencia}, e a
segunda tende a $\mu$. Portanto,
\begin{equation}\label{eq:subsequencia-final}
\frac{T_{k(n)}}{k(n)}\longrightarrow\mu
\qquad\text{quase certamente.}
\end{equation}

\subsection{Os Valores Intermediários}\label{subsec:intermediarios}

O limite quase certo garantido por \eqref{eq:subsequencia-final} fornece
a convergência apenas ao longo dos índices $k(n)$.
Para completar a prova falta controlar os índices $m$ que ficam entre dois valores
consecutivos da sequência $k(n)$. A hipótese $X_n\geqslant0$ será a chave para fazer o controle
da sequência ao longo desses índices, por meio de um argumento de monotonicidade.

Já que $Y_k\geqslant0$, então podemos afirmar que a sequência $T_m=Y_1+\cdots+Y_m$ é não decrescente em
$m$. Seja $m$ um inteiro positivo satisfazendo
\[
k(n)\leqslant m<k(n+1).
\]
Pela monotonicidade, $T_{k(n)}\leqslant T_m\leqslant T_{k(n+1)}$. Usando
também a não negatividade das variáveis aleatórias $T_n$'s e as desigualdades acima, obtemos
\[
\frac{T_{k(n)}}{k(n+1)}
\leqslant
\frac{T_m}{k(n+1)}
\leqslant
\frac{T_m}{m}
\leqslant
\frac{T_{k(n+1)}}{m}
\leqslant
\frac{T_{k(n+1)}}{k(n)},
\]
isto é,
\begin{equation}\label{eq:sanduiche}
\frac{T_{k(n)}}{k(n+1)}
\leqslant
\frac{T_m}{m}
\leqslant
\frac{T_{k(n+1)}}{k(n)}.
\end{equation}
Para usar \eqref{eq:sanduiche}, precisamos calcular os limites de seus dois
extremos. Em cada um deles, multiplicaremos e dividiremos pelo índice da
soma que aparece no numerador. Assim, faremos aparecer as médias cuja
convergência já conhecemos, acompanhadas de razões entre índices
consecutivos.

Comecemos por essas razões. Pela definição da parte inteira,
$0\leqslant\alpha^n-k(n)<1$. Dividindo por $\alpha^n$ e fazendo
$n\to\infty$, obtemos $k(n)/\alpha^n\to1$. Consequentemente,
\[
\frac{k(n+1)}{k(n)}
=
\frac{k(n+1)}{\alpha^{n+1}}\cdot\alpha\cdot\frac{\alpha^n}{k(n)}
\longrightarrow\alpha,
\]
e a razão inversa tende a $\alpha^{-1}$.

No que segue trabalhamos no evento de probabilidade $1$ em que vale
\eqref{eq:subsequencia-final}. Nesse evento,
\begin{gather*}
\frac{T_{k(n)}}{k(n+1)}
=
\frac{T_{k(n)}}{k(n)}\cdot\frac{k(n)}{k(n+1)}
\longrightarrow
\frac{\mu}{\alpha},
\\[0.4cm]
\frac{T_{k(n+1)}}{k(n)}
=
\frac{T_{k(n+1)}}{k(n+1)}\cdot\frac{k(n+1)}{k(n)}
\longrightarrow
\alpha\mu.
\end{gather*}

Na desigualdade \eqref{eq:sanduiche}, as duas cotas dependem apenas de $n$.
Assim, elas controlam de uma só vez todas as médias $T_m/m$ com
$k(n)\leqslant m<k(n+1)$. Vamos usar esse controle em cada intervalo para
obter uma estimativa válida para todos os índices suficientemente grandes.
A \Cref{fig:blocos-geometricos} ilustra como os intervalos se sucedem.

\begin{figure}[H]
\centering
\begin{tikzpicture}[x=1cm,y=1cm]
  % As distâncias são esquemáticas; importa a cobertura dos índices.
  \fill[teal!12] (0.4,-0.12) rectangle (3.2,0.12);
  \fill[teal!22] (3.2,-0.12) rectangle (7,0.12);
  \fill[teal!12] (7,-0.12) rectangle (11.5,0.12);
  \draw[->,thin] (0,0) -- (12.6,0);

  \foreach \x/\rotulo in {0.4/N,3.2/{N+1},7/{N+2},11.5/{N+3}} {
    \draw[thin] (\x,-0.18) -- (\x,0.18);
    \node[below=6pt,font=\small] at (\x,0) {$k(\rotulo)$};
  }
  \fill[teal] (5.1,0) circle (2pt);
  \node[above=4pt,teal!70!black] at (5.1,0) {$m$};

  \draw[muted,thin,decorate,decoration={brace,amplitude=4pt}]
    (0.4,0.65) -- (3.2,0.65)
    node[midway,above=6pt,font=\small] {$n=N$};
  \draw[muted,thin,decorate,decoration={brace,amplitude=4pt}]
    (3.2,0.65) -- (7,0.65)
    node[midway,above=6pt,font=\small] {$n=N+1$};
  \draw[muted,thin,decorate,decoration={brace,amplitude=4pt}]
    (7,0.65) -- (11.5,0.65)
    node[midway,above=6pt,font=\small] {$n=N+2$};
  \node[font=\small] at (12.1,0.65) {$\cdots$};
\end{tikzpicture}
\caption{Os intervalos $\left[k(n),k(n+1)\right)$, com $n\geqslant N$,
cobrem todos os inteiros $m\geqslant k(N)$. Em cada intervalo, a desigualdade
\eqref{eq:sanduiche} controla todas as médias $T_m/m$ pelas mesmas duas
cotas. O desenho é esquemático.}
\label{fig:blocos-geometricos}
\end{figure}

Fixemos $\delta>0$. Pelos dois limites que acabamos de obter, existe $N$
tal que, para todo $n\geqslant N$,
\[
\frac{\mu}{\alpha}-\delta
\leqslant\frac{T_{k(n)}}{k(n+1)}
\qquad\text{e}\qquad
\frac{T_{k(n+1)}}{k(n)}\leqslant\alpha\mu+\delta.
\]
Combinando essas cotas com \eqref{eq:sanduiche}, obtemos
\[
\frac{\mu}{\alpha}-\delta
\leqslant\frac{T_m}{m}\leqslant\alpha\mu+\delta
\qquad\text{se }k(n)\leqslant m<k(n+1),\quad n\geqslant N.
\]
Como $k(n)\to\infty$, esses intervalos cobrem todos os inteiros
$m\geqslant k(N)$. Portanto, a estimativa acima vale para todo
$m\geqslant k(N)$.

Uma vez que $\delta>0$ é arbitrário, o limite inferior das médias é pelo
menos $\mu/\alpha$, e o limite superior é no máximo $\alpha\mu$. Concluímos,
no evento de probabilidade $1$ em que estamos trabalhando, que
\begin{equation}\label{eq:limites-alfa}
\frac{\mu}{\alpha}
\leqslant
\liminf_{m\to\infty}\frac{T_m}{m}
\leqslant
\limsup_{m\to\infty}\frac{T_m}{m}
\leqslant
\alpha\mu.
\end{equation}

Para concluir que o limite é $\mu$, precisamos aproximar as duas cotas de
\eqref{eq:limites-alfa} desse mesmo valor. Isso sugere escolher $\alpha$
cada vez mais próximo de $1$, o que podemos fazer porque o argumento vale
para qualquer $\alpha>1$.

Precisamos, porém, que as desigualdades usadas nessa passagem valham no
mesmo evento de probabilidade $1$. Para cada $\alpha$ fixado, obtivemos um
evento com essa propriedade, que pode depender de $\alpha$. Basta então
escolher uma sequência enumerável de parâmetros que tenda a $1$:
\[
\alpha_j=1+\frac{1}{j},\qquad j\geqslant1.
\]
A interseção dos eventos de probabilidade $1$ correspondentes a esses
parâmetros também tem probabilidade $1$. Nessa interseção,
\eqref{eq:limites-alfa} vale simultaneamente com $\alpha=\alpha_j$ para
todo $j\geqslant1$. Podemos, portanto, fazer $j\to\infty$ nas duas cotas.
Como $\mu/\alpha_j\to\mu$ e $\alpha_j\mu\to\mu$, obtemos
\[
\mu
\leqslant
\liminf_{m\to\infty}\frac{T_m}{m}
\leqslant
\limsup_{m\to\infty}\frac{T_m}{m}
\leqslant
\mu.
\]
Os limites inferior e superior são, assim, iguais a $\mu$, o que prova que
$T_m/m\to\mu$ quase certamente. Pelo \Cref{lem:truncamento}, as médias
originais $S_m/m$ também convergem quase certamente para $\mu$. Isso conclui
a prova do \Cref{thm:lei-forte} para variáveis não negativas; a redução da
\Cref{subsec:reducao} fornece o caso geral.

\subsection{Observações sobre a Prova}\label{subsec:observacoes-prova}

A prova está completa, mas vale destacar três pontos que costumam passar
despercebidos na primeira leitura.

O primeiro é o papel de cada hipótese. A hipótese de distribuições iguais é
usada no lema de truncamento, para escrever
$\mathbb{P}(|X_k|>k)=\mathbb{P}(|X_1|>k)$, na estimativa dos segundos
momentos, para substituir $\mathbb{P}(|X_k|>y)$ pela cauda de $X_1$, e na
identificação de $\mathbb{E}[Y_k]$ com a esperança do truncamento
correspondente de $X_1$. Junto com a fórmula de integração da cauda e a cota
do \Cref{lem:harmonica}, isso produz a cota $4\mathbb{E}[|X_1|]$ no
\Cref{lem:momento}. A hipótese de independência dois a dois é usada apenas
para que a variância da soma seja a soma das variâncias; não é preciso mais do
que isso. A hipótese
$\mathbb{E}[|X_1|]<\infty$ é usada no truncamento, para garantir a
somabilidade de \eqref{eq:serie-truncamento}, e na estimativa de momentos, para
garantir a finitude do lado direito do \Cref{lem:momento}. Ela também garante
que as médias truncadas convergem para um centro finito em \eqref{eq:cesaro}.

O segundo é o papel do parâmetro $\alpha$. Ele é arbitrário em $(1,\infty)$ e o
argumento funciona para qualquer valor fixado. Como o conjunto nulo pode
depender de $\alpha$, na conclusão intersectamos os conjuntos de probabilidade
$1$ apenas para a sequência enumerável $\alpha_j=1+1/j$. Fazendo $j\to\infty$,
identificamos o limite. Se tivéssemos fixado $\alpha=2$, por exemplo,
obteríamos apenas
\[
\frac{\mu}{2}\leqslant\liminf\frac{T_m}{m}\leqslant\limsup\frac{T_m}{m}\leqslant2\mu,
\]
que, quando $\mu>0$, ainda não identifica o limite.

O terceiro é o papel da não negatividade. Ela garante tanto a monotonicidade
de $T_m$ quanto os sinais necessários às comparações dos denominadores em
\eqref{eq:sanduiche}. Sem $Y_k\geqslant0$, essas desigualdades podem falhar e não teríamos o
controle dos valores intermediários usado neste argumento.

%----------------------------------------------------------------------
\section{A Recíproca: A Integrabilidade é Necessária}\label{sec:reciproca}

A hipótese $\mathbb{E}[|X_1|]<\infty$ não é uma conveniência técnica. Nesta
seção, mostramos que, sem ela, a conclusão da lei forte falha para sequências
independentes e identicamente distribuídas. O argumento é o da Seção 2.3 de
\cite{MR3930614}.

\begin{theorem}[Recíproca da lei forte]\label{thm:reciproca}
Sejam $X_1,X_2,\ldots$ variáveis aleatórias independentes e identicamente
distribuídas com $\mathbb{E}[|X_1|]=\infty$. Então
\[
\mathbb{P}\bigl(|X_n|\geqslant n\ \text{infinitas vezes}\bigr)=1.
\]
Consequentemente, se $S_n=X_1+\cdots+X_n$,
\[
\mathbb{P}\Bigl(\lim_{n\to\infty}\frac{S_n}{n}\ \text{existe e é finito}\Bigr)=0.
\]
\end{theorem}

\begin{proof}
Aplicando a fórmula \eqref{eq:esperanca-cauda} e a comparação
\eqref{eq:comparacao-cauda} a $W=|X_1|$, obtemos
\begin{equation}\label{eq:cauda-serie}
\mathbb{E}[|X_1|]
=
\int_0^{\infty}\mathbb{P}(|X_1|>t)\,dt
\leqslant
\sum_{n=0}^{\infty}\mathbb{P}(|X_1|>n).
\end{equation}
Como a esperança é infinita, a série do lado direito diverge. Sendo as
variáveis identicamente distribuídas,
\[
\sum_{n=1}^{\infty}\mathbb{P}(|X_n|\geqslant n)
\geqslant
\sum_{n=1}^{\infty}\mathbb{P}(|X_n|>n)
=
\sum_{n=1}^{\infty}\mathbb{P}(|X_1|>n)
=
\infty.
\]
Os eventos $\{|X_n|\geqslant n\}$ são independentes, logo o segundo lema de
Borel-Cantelli fornece
\begin{equation}\label{eq:io-recorde}
\mathbb{P}\bigl(|X_n|\geqslant n\ \text{infinitas vezes}\bigr)=1.
\end{equation}

Para concluir, mostraremos que a convergência de $S_n/n$ a um limite finito
obrigaria $X_n/n\to0$, em contradição com \eqref{eq:io-recorde}.
Suponha que exista $\omega$ tal que $S_n(\omega)/n$ convirja para um número
real $\ell$. Da identidade
\[
\frac{S_n}{n}-\frac{S_{n+1}}{n+1}
=
\frac{S_n}{n(n+1)}-\frac{X_{n+1}}{n+1},
\]
obtemos
\[
\frac{X_{n+1}}{n+1}
=
\frac{S_n}{n(n+1)}-\Bigl(\frac{S_n}{n}-\frac{S_{n+1}}{n+1}\Bigr).
\]
O primeiro termo do lado direito vale
$\frac{1}{n+1}\cdot\frac{S_n}{n}$, que tende a zero porque $S_n/n\to\ell$ e
$\ell$ é finito. O segundo tende a $\ell-\ell=0$. Logo
\begin{equation}\label{eq:xn-sobre-n}
\frac{|X_{n+1}(\omega)|}{n+1}\longrightarrow0.
\end{equation}
Por outro lado, \eqref{eq:xn-sobre-n} é equivalente, após reindexação, a
$|X_n(\omega)|/n\to0$. Mas o evento
$\{|X_n|\geqslant n\ \text{infinitas vezes}\}$ significa exatamente que
$|X_n(\omega)|/n\geqslant1$ para infinitos índices $n$, o que contradiz a
convergência para zero. Portanto o conjunto em que $S_n/n$
converge para um limite finito é disjunto de
$\{|X_n|\geqslant n\ \text{infinitas vezes}\}$. Por \eqref{eq:io-recorde}, esse
último conjunto tem probabilidade $1$, e a conclusão segue.
\end{proof}

A ideia do argumento é identificar uma condição necessária para a
convergência das médias: se $S_n/n$ converge para um limite finito, então
$X_n/n\to0$. Sem integrabilidade, porém, quase certamente
$|X_n|\geqslant n$ para infinitos índices $n$. Assim, $|X_n|/n$ é pelo menos
$1$ infinitas vezes, o que impede essa razão de tender a zero. Portanto, as
médias também não podem convergir para um limite finito.

\begin{example}[A distribuição de Cauchy]\label{ex:cauchy}
Sejam $X_1,X_2,\ldots$ independentes com a distribuição de Cauchy padrão, cuja
densidade é
\[
f(x)=\frac{1}{\pi(1+x^2)},
\qquad x\in\mathbb{R}.
\]
Então
\[
\mathbb{E}[|X_1|]
=
\int_{-\infty}^{\infty}\frac{|x|}{\pi(1+x^2)}\,dx
=
\frac{2}{\pi}\int_0^{\infty}\frac{x}{1+x^2}\,dx
=
\infty,
\]
de modo que o \Cref{thm:reciproca} se aplica: $S_n/n$ não converge quase
certamente.

Por outro lado, a média de variáveis de Cauchy independentes também é de
Cauchy. Este fato será provado mais a frente quando introduzirmos a função característica.
De qualquer forma usando este resultado podemos afirmar que $S_n/n$ tem
exatamente a mesma distribuição para todo $n$ e, em particular,
$S_n/n\Rightarrow X_1$. O exemplo mostra que a convergência em distribuição
pode valer trivialmente enquanto a convergência quase certa falha, e que a
hipótese de integrabilidade é indispensável para a lei forte.
\end{example}

No exemplo da Cauchy, a média também não converge em probabilidade para
nenhuma constante $c\in\mathbb{R}$ já que para todo $\varepsilon>0$,
\[
\mathbb{P}(|S_n/n-c|\geqslant\varepsilon)
=\mathbb{P}(|X_1-c|\geqslant\varepsilon)>0,
\]
independentemente de $n$. Isso não contradiz
a lei fraca, porque a lei fraca na forma demonstrada exige variância finita, e
a Cauchy não tem nem esperança. No caso independente e identicamente
distribuído, a integrabilidade é suficiente, mas não necessária, para a lei
fraca. Já o \Cref{thm:reciproca} mostra que ela é necessária para a
convergência quase certa das médias a um limite finito.

%----------------------------------------------------------------------
\section{Aplicações}\label{sec:aplicacoes}

\subsection{Números Normais}\label{subsec:numeros-normais}

O \Cref{thm:borel-reais} mostra que, para quase todo número real, cada
dígito decimal tem frequência assintótica $1/10$. Podemos perguntar também
pelas frequências de pares, trincas e blocos mais longos. Por exemplo, será
que os blocos $00$ e $37$ têm, cada um, frequência $1/100$?

O equilíbrio dos dígitos isolados não basta para responder a essa pergunta.
Considere o número
\[
x=0{,}0123456789\,0123456789\,0123456789\ldots.
\]
Cada dígito aparece uma vez a cada dez posições, mas o bloco $00$ nunca
aparece. Esse número é simplesmente normal na base $10$, mas não é normal.
A normalidade, definida na \Cref{def:normal}, exige que \emph{toda} palavra
finita $u$ de comprimento $\ell$ tenha frequência $10^{-\ell}$. Vamos mostrar
que quase todo número real satisfaz essa propriedade.

\subsubsection{Contar as ocorrências de uma palavra}

Retomemos o espaço $[0,1)$ com a probabilidade de Lebesgue $\lambda$ e os
dígitos $d_j(x)$ da expansão decimal de $x$, com a convenção da
\Cref{subsec:digitos-reais}. Pelo \Cref{lem:digitos-uniformes}, essas
variáveis são independentes e uniformes em $\{0,\ldots,9\}$.

Fixemos uma palavra $u=u_1\ldots u_\ell$ de comprimento $\ell\geqslant1$.
Para registrar se ela aparece a partir da posição $j$, definimos
\[
Z_j^{(u)}(x)
=\mathds{1}_{\{d_j=u_1,\ldots,d_{j+\ell-1}=u_\ell\}}(x).
\]
Assim, $Z_j^{(u)}(x)=1$ quando o bloco que começa em $j$ é $u$, e vale $0$
caso contrário. A frequência de ocorrências entre as $n$ primeiras
posições de início é, portanto,
\[
F_n^{(u)}(x)=\frac{1}{n}\sum_{j=1}^{n}Z_j^{(u)}(x).
\]
Contamos todas as posições de início, inclusive as que produzem ocorrências
sobrepostas. Por exemplo, o bloco $00$ aparece duas vezes em $000$, com
inícios na primeira e na segunda posições.

Pela independência dos dígitos, a probabilidade de encontrar $u$ em qualquer
posição fixada é
\[
\lambda\bigl(Z_j^{(u)}=1\bigr)
=\prod_{i=1}^{\ell}\lambda\bigl(d_{j+i-1}=u_i\bigr)
=\left(\frac{1}{10}\right)^\ell
=10^{-\ell}.
\]
Queremos mostrar que a frequência observada $F_n^{(u)}$ converge quase
certamente para essa probabilidade.

\subsubsection{Separar os blocos para obter independência}

Uma tentativa seria aplicar a lei forte diretamente à sequência
$(Z_j^{(u)})$. A dificuldade é que blocos com posições iniciais próximas
usam alguns dos mesmos dígitos. Por exemplo, se $u=372$, a condição
$Z_1^{(u)}=1$ exige $d_2=7$, enquanto $Z_2^{(u)}=1$ exige $d_2=3$.
Esses dois eventos têm probabilidade positiva, mas não podem ocorrer
simultaneamente. Logo, as indicadoras não são independentes.

A sobreposição sugere como proceder: separar as posições iniciais de modo
que os blocos examinados usem posições disjuntas da expansão decimal. Para palavras
de comprimento $3$, os blocos que começam em $1,4,7,\ldots$ examinam,
respectivamente,
\[
(d_1,d_2,d_3),\qquad(d_4,d_5,d_6),\qquad(d_7,d_8,d_9),\qquad\ldots.
\]
Como os dígitos são independentes e nenhum deles é usado em dois desses
blocos, as indicadoras correspondentes são independentes. Podemos fazer
o mesmo com os inícios $2,5,8,\ldots$ e com os inícios $3,6,9,\ldots$.
Esses três grupos, reunidos, incluem todas as posições de início, como
mostra a \Cref{fig:blocos-normais}.

\begin{figure}[H]
\centering
\begin{tikzpicture}[x=0.88cm,y=1cm]
  % As colunas mantêm a mesma posição decimal em todas as linhas.
  \node[anchor=west,font=\small\bfseries] at (-2.3,0.85)
    {Testes em posições consecutivas};
  \foreach \y/\inicio in {0/1,-0.75/2} {
    \node[anchor=east,font=\small] at (-0.2,\y) {início $\inicio$};
    \foreach \j in {1,...,4} {
      \draw[muted!35,thin] ({\j-1},{\y-0.24}) rectangle (\j,{\y+0.24});
      \node[muted!65,font=\small] at ({\j-0.5},\y) {$d_{\j}$};
    }
    \fill[teal!16] ({\inicio-1},{\y-0.24})
      rectangle ({\inicio+2},{\y+0.24});
    \foreach \s in {0,1,2} {
      \pgfmathtruncatemacro{\j}{\inicio+\s}
      \draw[muted!40,thin] ({\j-1},{\y-0.24}) rectangle (\j,{\y+0.24});
      \node[font=\small] at ({\j-0.5},\y) {$d_{\j}$};
    }
    \draw[teal,thin] ({\inicio-0.97},{\y-0.27})
      rectangle ({\inicio+1.97},{\y+0.27});
  }
  \foreach \x in {1.5,2.5} {
    \draw[teal!70,dashed,thin] (\x,-0.28) -- (\x,-0.47);
  }
  \node[anchor=west,align=left,font=\small,text width=5.3cm]
    at (5,-0.35)
    {Os dois testes usam os dígitos $d_2$ e $d_3$.\\
     A sobreposição pode produzir dependência.};

  \node[anchor=west,font=\small\bfseries] at (-2.3,-1.7)
    {Testes separados em três grupos};
  \foreach \r in {1,2,3} {
    \pgfmathsetmacro{\y}{-2.45-0.85*(\r-1)}
    \pgfmathtruncatemacro{\segundo}{\r+3}
    \pgfmathtruncatemacro{\terceiro}{\r+6}
    \node[anchor=east,font=\small] at (-0.2,\y)
      {inícios $\r,\segundo,\terceiro,\ldots$};
    \foreach \j in {1,...,11} {
      \node[muted!55,font=\small] at ({\j-0.5},\y) {$d_{\j}$};
    }
    \foreach \k in {0,1,2} {
      \pgfmathtruncatemacro{\inicio}{\r+3*\k}
      \fill[teal!16] ({\inicio-1},{\y-0.24})
        rectangle ({\inicio+2},{\y+0.24});
      \foreach \s in {0,1,2} {
        \pgfmathtruncatemacro{\j}{\inicio+\s}
        \draw[muted!40,thin] ({\j-1},{\y-0.24}) rectangle (\j,{\y+0.24});
        \node[font=\small] at ({\j-0.5},\y) {$d_{\j}$};
      }
      \draw[teal,thin] ({\inicio-0.97},{\y-0.27})
        rectangle ({\inicio+1.97},{\y+0.27});
    }
    \node[font=\small] at (11.5,\y) {$\cdots$};
  }
\end{tikzpicture}
\caption{Cada retângulo destacado representa um bloco de três dígitos
examinado na procura de uma palavra fixada. Acima, os testes compartilham
dígitos. Abaixo, os blocos de uma mesma linha usam posições disjuntas,
o que garante a independência das indicadoras nessa linha. As três linhas
incluem, juntas, todas as posições de início.}
\label{fig:blocos-normais}
\end{figure}

Para uma palavra de comprimento $\ell$, fazemos a mesma separação em
$\ell$ grupos. Fixado $r\in\{1,\ldots,\ell\}$, consideramos as posições
\[
r,\quad r+\ell,\quad r+2\ell,\quad\ldots.
\]
As indicadoras $Z_r^{(u)},Z_{r+\ell}^{(u)},Z_{r+2\ell}^{(u)},\ldots$ são
independentes, têm a mesma distribuição de Bernoulli com parâmetro
$10^{-\ell}$ e são limitadas por $1$. Pelo \Cref{cor:limitadas},
\[
\frac{1}{q}\sum_{k=0}^{q-1}Z_{r+k\ell}^{(u)}
\longrightarrow10^{-\ell}
\qquad\text{quase certamente.}
\]
Aplicamos esse argumento a cada um dos $\ell$ grupos. A interseção dos
$\ell$ eventos de probabilidade $1$ também tem probabilidade $1$; nela,
todas essas frequências convergem simultaneamente.

\subsubsection{Reunir as frequências dos grupos}

\begin{samepage}
Vamos primeiro contar as ocorrências com início nas primeiras $q\ell$
posições. Entre elas, há exatamente $q$ posições de cada grupo. Portanto,
a frequência total é a média aritmética das $\ell$ frequências que acabamos
de estudar:
\[
F_{q\ell}^{(u)}
=\frac{1}{\ell}\sum_{r=1}^{\ell}
\left(\frac{1}{q}\sum_{k=0}^{q-1}Z_{r+k\ell}^{(u)}\right)
\longrightarrow
\frac{1}{\ell}\sum_{r=1}^{\ell}10^{-\ell}
=10^{-\ell}.
\]
\end{samepage}
Essa identidade explica como a separação em grupos preserva a contagem de
todas as ocorrências, inclusive as sobrepostas.

Para tratar todos os $n\geqslant\ell$, escrevemos
$n=q\ell+s$, com $0\leqslant s<\ell$. Além das primeiras $q\ell$ posições,
restam apenas $s$ posições a examinar. Assim,
\[
F_n^{(u)}
=\frac{q\ell}{n}\,F_{q\ell}^{(u)}
+\frac{1}{n}\sum_{j=q\ell+1}^{n}Z_j^{(u)}.
\]
A última parcela está entre $0$ e $(\ell-1)/n$, pois cada indicadora vale
no máximo $1$. Como $n-q\ell=s$ permanece limitado, temos $q\to\infty$ e
$q\ell/n\to1$ quando $n\to\infty$. Obtemos, portanto,
\begin{equation}\label{eq:frequencia-palavra}
F_n^{(u)}\longrightarrow10^{-\ell}
\qquad\text{quase certamente.}
\end{equation}

Se exigirmos que cada ocorrência esteja inteiramente contida nos primeiros
$n$ dígitos, como no \Cref{exercise:palavras}, basta excluir os inícios
posteriores à posição $n-\ell+1$. São no máximo $\ell-1$ ocorrências a
menos, de modo que a diferença entre as duas frequências é no máximo
$(\ell-1)/n$ e tende a zero. As duas maneiras de contar fornecem o mesmo
limite.

\subsubsection{Todas as palavras e todas as bases}

Até aqui, fixamos uma palavra $u$. Para cada comprimento $\ell$, existem
apenas $10^\ell$ palavras, de modo que a família de todas as palavras
finitas é enumerável. Para cada uma delas, \eqref{eq:frequencia-palavra}
falha apenas em um conjunto de medida zero. A união desses conjuntos
excepcionais também tem medida zero. Fora dela, todas as palavras têm,
simultaneamente, a frequência exigida pela normalidade.

Concluímos que quase todo $x\in[0,1)$ é normal na base $10$. Como na
\Cref{subsec:digitos-reais}, o resultado se estende a $\mathbb{R}$ por
translações inteiras: os dígitos considerados são os da parte fracionária,
e as translações do conjunto excepcional ainda formam uma união
enumerável de conjuntos de medida zero.

Por fim, a escolha da base $10$ não foi essencial. Em uma base inteira
$b\geqslant2$, dividimos o intervalo em $b$ partes iguais a cada etapa.
Os dígitos são então independentes e uniformes em $\{0,\ldots,b-1\}$, e
a probabilidade de uma palavra de comprimento $\ell$ é $b^{-\ell}$.
O mesmo argumento mostra que quase todo número real é normal nessa base.

Um número é chamado \emph{completamente normal} quando é normal em toda
base inteira $b\geqslant2$. Para cada base, acabamos de obter um conjunto
excepcional de medida zero. Como há apenas uma quantidade enumerável de
bases inteiras, a união desses conjuntos também tem medida zero.
Fora dela, a normalidade vale em todas as bases simultaneamente. Portanto,
quase todo número real é completamente normal.

\subsection{Integração de Monte Carlo}\label{subsec:monte-carlo}

O método de Monte Carlo para calcular uma integral se apoia na lei forte. Seja
$f:[0,1]\to\mathbb{R}$ uma função Borel-mensurável com
$\int_0^1|f(x)|\,dx<\infty$, e sejam $U_1,U_2,\ldots$
independentes e uniformes em $[0,1]$. Então as variáveis $f(U_i)$ são
independentes e identicamente distribuídas, com
\[
\mathbb{E}\bigl[|f(U_1)|\bigr]=\int_0^1|f(x)|\,dx<\infty.
\]
Pelo \Cref{thm:lei-forte},
\begin{equation}\label{eq:monte-carlo}
\frac{1}{n}\sum_{i=1}^{n}f(U_i)\longrightarrow\int_0^1f(x)\,dx
\qquad\text{quase certamente.}
\end{equation}
A integral é o valor esperado de $f(U_1)$, e a média empírica é o estimador
obtido por amostragem. A lei forte garante que o estimador converge ao longo de
quase toda realização da sequência, e não apenas em probabilidade. Se,
adicionalmente, $\int_0^1 f(x)^2\,dx<\infty$, a desigualdade de Chebyshev
fornece uma cota para a probabilidade de erro em cada $n$ finito. Sob as
hipóteses usuais do teorema central do limite, aqui satisfeitas quando
$f(U_1)$ tem variância finita e positiva, obtemos também uma aproximação
normal assintótica para o erro; essa aproximação não é uma cota para $n$
finito.

Tomando $f=\mathds{1}_{[k/10,(k+1)/10)}$, a indicadora aplicada a cada $U_i$
registra se o primeiro dígito decimal da $i$-ésima amostra é $k$. A média
empírica é, portanto, a frequência desse dígito entre as amostras, e a integral
vale $1/10$. Esses indicadores têm a mesma distribuição conjunta que os
indicadores dos dígitos sucessivos de um único número uniforme em $[0,1]$:
ambas as sequências são formadas por variáveis independentes com distribuição
de Bernoulli com parâmetro $1/10$.

\subsection{A Função de Distribuição Empírica}\label{subsec:empirica}

Sejam $X_1,X_2,\ldots$ independentes e identicamente distribuídas com função de
distribuição $F$. Para cada $x\in\mathbb{R}$, a variável
$\mathds{1}_{\{X_i\leqslant x\}}$ é de Bernoulli com parâmetro $F(x)$, e as
variáveis são independentes. Pelo \Cref{thm:lei-forte},
\begin{equation}\label{eq:empirica-pontual}
F_n(x)\equiv\frac{1}{n}\sum_{i=1}^{n}\mathds{1}_{\{X_i\leqslant x\}}
\longrightarrow F(x)
\qquad\text{quase certamente,}
\end{equation}
para cada $x$ fixado. A função $F_n$ é a \emph{função de distribuição empírica}
e $F_n(x)$ é a frequência observada de valores menores ou iguais a $x$.

A afirmação acima é pontual: para cada $x$ fixado, vale quase certamente, mas
isso, por si só, não fornece um único evento de probabilidade $1$ que sirva para
todos os $x$. Primeiro suponha que $F$ seja contínua. Como $\mathbb{Q}$ é
enumerável, há um evento de probabilidade $1$ no qual $F_n(q)\to F(q)$ para
todo racional $q$. Fixado $\varepsilon>0$, escolha uma malha finita racional
$q_1<\cdots<q_m$ que cubra a parte central da distribuição, com as massas
entre pontos consecutivos menores que $\varepsilon$ e com as duas caudas
menores que $\varepsilon$. Para construir essa malha, escolhemos primeiro
extremos racionais com caudas pequenas, usando os limites de $F$ em
$-\infty$ e $+\infty$. A continuidade uniforme de $F$ no intervalo compacto
entre esses extremos permite então subdividi-lo em uma malha racional com
as variações desejadas. Para $q_j\leqslant x\leqslant q_{j+1}$, a
monotonicidade dá, em ambas as direções,
\[
|F_n(x)-F(x)|\leqslant
\max_{1\leqslant j\leqslant m}|F_n(q_j)-F(q_j)|+\varepsilon,
\]
e os mesmos controles valem nas caudas. Como a malha é finita, o máximo dos
erros nos seus pontos tende a zero. Aplicando o argumento a uma sequência
enumerável de $\varepsilon$ que tende a zero, obtemos
\[
\sup_{x\in\mathbb{R}}\bigl|F_n(x)-F(x)\bigr|\longrightarrow0
\qquad\text{quase certamente.}
\]
Esse é o teorema de Glivenko-Cantelli no caso contínuo. Se $F$ tem átomos, a
mesma intercalação por racionais garante a convergência nos pontos de
continuidade, mas não garante a convergência em todos os saltos, que podem
ocorrer em pontos irracionais. Nos saltos racionais, a convergência já está
incluída no evento escolhido. Os átomos formam um
conjunto enumerável, e a convergência em cada posição atômica também vale quase
certamente pela lei forte, assim como a convergência dos limites à esquerda
$F_n(a-)$ para $F(a-)$ em cada átomo $a$. Para obter a extensão uniforme,
separam-se os finitos átomos de massa maior que uma tolerância fixada e refinam-se os
intervalos entre eles até que as variações restantes de $F$ sejam pequenas; a
última parte do \Cref{exercise:glivenko} detalha essa ideia.

\subsection{Distribuições de Cauda Pesada: A Família de Pareto}\label{subsec:pareto}

Um dos aspectos mais marcantes do \Cref{thm:lei-forte} é que ele exige apenas a
finitude do primeiro momento ($\mathbb{E}[|X_1|]<\infty$), dispensando
completamente a existência de variância ou de momentos de ordem superior. Para
compreender o alcance dessa generalidade, é instrutivo examinar situações em que
a média teórica existe, mas as flutuações ao redor dela são dominadas por caudas
pesadas.

Um modelo clássico com esse comportamento é a \emph{distribuição de Pareto}.
Sejam $\alpha>0$ e $x_0>0$. Uma variável aleatória contínua $X$ tem distribuição
de Pareto com parâmetro de forma $\alpha$ e escala $x_0$ se sua densidade é dada
por
\begin{equation}\label{eq:densidade-pareto}
f(x)=
\begin{cases}
\dfrac{\alpha x_0^\alpha}{x^{\alpha+1}}, & \text{para } x\geqslant x_0,\\[6pt]
0, & \text{para } x<x_0.
\end{cases}
\end{equation}
Para fixar ideias e simplificar a álgebra, fixemos a escala $x_0=1$. A
probabilidade de cauda decai como uma lei de potência:
\[
\mathbb{P}(X>t)=\int_t^{\infty}\frac{\alpha}{x^{\alpha+1}}\,dx=\frac{1}{t^\alpha},
\qquad t\geqslant1.
\]
Calculando os momentos de ordem $p>0$, temos
\[
\mathbb{E}[X^p]=\int_1^{\infty}x^p\frac{\alpha}{x^{\alpha+1}}\,dx
=\alpha\int_1^{\infty}x^{p-\alpha-1}\,dx,
\]
integral que converge se, e somente se, $p<\alpha$.

Consideremos agora o regime $1<\alpha\leqslant2$, tomando por exemplo
$\alpha=3/2$. Nesse caso:
\begin{enumerate}
\item[(i)] O primeiro momento é finito:
\[
\mathbb{E}[X_1]=\frac{\alpha}{\alpha-1}=\frac{3/2}{1/2}=3<\infty;
\]
\item[(ii)] O segundo momento diverge:
\[
\mathbb{E}[X_1^2]=\frac{3}{2}\int_1^{\infty}x^{2-5/2}\,dx
=\frac{3}{2}\int_1^{\infty}\frac{1}{\sqrt{x}}\,dx=\infty,
\]
e, consequentemente, $\operatorname{Var}(X_1)=\infty$ e $\mathbb{E}[X_1^4]=\infty$.
\end{enumerate}

Essa classe de distribuições desempenha um papel fundamental em finanças,
ciência de redes e teoria do risco, modelando fenômenos sujeitos a eventos
extremos (como grandes sinistros em companhias de seguros, distribuições de renda
ou volumes de tráfego de dados na internet).

Suponha agora que $X_1,X_2,\ldots$ sejam variáveis aleatórias independentes com
essa distribuição e denote por $S_n=X_1+\cdots+X_n$ a soma acumulada. O que
podemos afirmar sobre o comportamento assintótico da média empírica $S_n/n$?
\begin{itemize}
\item \textbf{A abordagem pelo quarto momento (\Cref{thm:quarto-momento}) não se
aplica:} como $\mathbb{E}[X_1^4]=\infty$, a estimativa via $(S_n-n\mu)^4$ sequer
está bem definida;
\item \textbf{A abordagem de Chebyshev e a Lei Fraca elementar falham:}
como $\operatorname{Var}(X_1)=\infty$, a desigualdade de Chebyshev não fornece
nenhum controle útil para $\mathbb{P}(|S_n/n-3|\geqslant\varepsilon)$;
\item \textbf{A Lei Forte de Etemadi (\Cref{thm:lei-forte}) aplica-se
plenamente:} como as variáveis são independentes e $\mathbb{E}[|X_1|]=3<\infty$,
o teorema garante que
\begin{equation}\label{eq:convergencia-pareto}
\frac{S_n}{n}\longrightarrow 3
\qquad\text{quase certamente.}
\end{equation}
\end{itemize}
Esse exemplo ilustra de forma nítida o poder da demonstração de Etemadi: a
convergência quase certa prevalece mesmo quando as observações individuais
produzem flutuações e saltos de variância infinita.

%----------------------------------------------------------------------
\section{Exercícios Propostos}\label{sec:exercicios}

\begin{exercise}\label{exercise:cauda-serie}
Seja $X$ uma variável aleatória. Dê uma demonstração alternativa da comparação
\eqref{eq:comparacao-cauda}, aplicada a $W=|X|$:
\[
\sum_{k=1}^{\infty}\mathbb{P}(|X|>k)
\leqslant
\mathbb{E}[|X|]
\leqslant
\sum_{k=0}^{\infty}\mathbb{P}(|X|>k),
\]
e conclua que a série $\sum_k\mathbb{P}(|X|>k)$ converge se, e somente se,
$\mathbb{E}[|X|]<\infty$. \emph{Sugestão:} compare pontualmente $|X|$ com as
somas $\sum_{k=1}^{\infty}\mathds{1}_{\{|X|>k\}}$ e
$\sum_{k=0}^{\infty}\mathds{1}_{\{|X|>k\}}$ e tome esperanças.
\end{exercise}

\begin{exercise}\label{exercise:bernoulli-l4}
Sejam $X_1,X_2,\ldots$ independentes com $\mathbb{P}(X_i=1)=p$ e
$\mathbb{P}(X_i=0)=1-p$, e $S_n=X_1+\cdots+X_n$.
\begin{enumerate}
\item[(a)] Calcule $\mathbb{E}[(S_n-np)^4]$ em função de $n$ e de $p$ e mostre
que ela é $O(n^2)$. Mais precisamente, para $0<p<1$, mostre que é da ordem de
$n^2$; trate também os casos degenerados $p=0$ e $p=1$.
\item[(b)] Aplique a desigualdade de Markov a $|S_n-np|^4$ para mostrar que
$\mathbb{P}(|S_n/n-p|\geqslant\varepsilon)\leqslant C/(n^2\varepsilon^4)$ e
conclua a convergência quase certa.
\end{enumerate}
\end{exercise}

\begin{exercise}\label{exercise:palavras}
No espaço $\Omega=\{0,1,\ldots,9\}^{\mathbb{N}}$ com a medida produto uniforme,
seja $u=u_1\ldots u_\ell$ uma palavra de comprimento $\ell\geqslant1$. Defina
\[
N_n(u)=\#\{1\leqslant j\leqslant n-\ell+1:X_j=u_1,\ldots,X_{j+\ell-1}=u_\ell\}.
\]
Mostre que $N_n(u)/n\to10^{-\ell}$ quase certamente. \emph{Sugestão:} considere
primeiro as posições $j=1,1+\ell,1+2\ell,\ldots$ e depois intercale as $\ell$
progressões aritméticas de razão $\ell$.
\end{exercise}

\begin{exercise}\label{exercise:monte-carlo}
Seja $f:[0,1]\to\mathbb{R}$ Borel-mensurável, com
$\int_0^1|f(x)|\,dx<\infty$, e sejam $U_1,U_2,\ldots$ independentes
e uniformes em $[0,1]$.
\begin{enumerate}
\item[(a)] Mostre que
$\frac{1}{n}\sum_{i=1}^{n}f(U_i)\to\int_0^1f(x)\,dx$ quase certamente.
\item[(b)] Suponha que $|f(x)|\leqslant M<\infty$ para todo $x\in[0,1]$.
Para $\varepsilon>0$, use a desigualdade de Chebyshev para estimar
$\mathbb{P}(|\frac{1}{n}\sum_{i=1}^{n}f(U_i)-\int_0^1f|\geqslant\varepsilon)$
em função de $n$, $\varepsilon$ e $M$.
\end{enumerate}
\end{exercise}

\begin{exercise}\label{exercise:media-infinita}
Sejam $X_1,X_2,\ldots$ independentes, identicamente distribuídas e não
negativas, com $\mathbb{E}[X_1]=\infty$, e defina
$S_n=\sum_{i=1}^{n}X_i$. Mostre que
\[
\frac{S_n}{n}\longrightarrow\infty
\qquad\text{quase certamente.}
\]
\emph{Sugestão:} para $M>0$, aplique a lei forte às variáveis truncadas
$X_i^M=\min\{X_i,M\}$ e use a convergência monótona para mostrar que
$\mathbb{E}[X_1^M]\to\infty$.
\end{exercise}

\begin{exercise}\label{exercise:pareto}
Sejam $X_1,X_2,\ldots$ independentes e identicamente distribuídas com
densidade $f(x)=x^{-2}$ para $x\geqslant1$ e $f(x)=0$ para $x<1$. Defina
$S_n=\sum_{i=1}^{n}X_i$.
\begin{enumerate}
\item[(a)] Mostre que $\mathbb{E}[|X_1|]=\infty$ e que
$\mathbb{P}(X_1>t)=1/t$ para $t\geqslant1$.
\item[(b)] Conclua, usando o exercício anterior, que $S_n/n$ não converge quase
certamente para um limite finito.
\item[(c)] Mostre que $\mathbb{P}(|X_n|\geqslant n\ \text{infinitas vezes})=1$,
usando diretamente o segundo lema de Borel-Cantelli.
\end{enumerate}
\end{exercise}

\begin{exercise}\label{exercise:momento-quatro}
Mostre que existe uma variável aleatória $X$ com $\mathbb{E}[|X|]<\infty$ e
$\mathbb{E}[X^4]=\infty$. Conclua que a hipótese do
\Cref{thm:quarto-momento} é estritamente mais forte que a do
\Cref{thm:lei-forte}.
\end{exercise}

\begin{starredexercise}\label{exercise:glivenko}
Sejam $X_1,X_2,\ldots$ independentes e identicamente distribuídas com função de
distribuição $F$ contínua, e seja $F_n$ a função de distribuição empírica.
\begin{enumerate}
\item[(a)] Mostre que, para cada $x$, $F_n(x)\to F(x)$ quase certamente.
\item[(b)] Mostre que, quase certamente, $F_n(q)\to F(q)$ simultaneamente para
todo racional $q$.
\item[(c)] Para cada $\varepsilon>0$, escolha um número finito de pontos
racionais $x_1<\cdots<x_k$ que cubra a parte central da distribuição, de modo
que as variações de $F$ entre pontos consecutivos e nas duas caudas sejam
menores que $\varepsilon$. Justifique a existência da malha escolhendo
primeiro extremos racionais com caudas pequenas e usando a continuidade
uniforme de $F$ no compacto entre eles. Considere também os extremos
$x_0=-\infty$ e $x_{k+1}=+\infty$, convencionando $F_n(-\infty)=F(-\infty)=0$ e
$F_n(+\infty)=F(+\infty)=1$. Usando a monotonicidade de $F_n$ e de $F$, mostre
que
\[
\sup_{x\in\mathbb{R}}|F_n(x)-F(x)|\leqslant
\max_{0\leqslant j\leqslant k+1}|F_n(x_j)-F(x_j)|+\varepsilon
\]
e conclua o teorema de Glivenko-Cantelli no caso contínuo.
\item[(d)] Agora suponha que $F$ possa ter átomos. Explique por que a malha de
(c), com variações menores que $\varepsilon$, pode não existir. Mostre que a
convergência vale simultaneamente em cada ponto atômico (há no máximo
enumeravelmente muitos) e também para os limites à esquerda nesses pontos,
observando que
\[
F_n(a-)=\frac{1}{n}\sum_{i=1}^{n}\mathds{1}_{\{X_i<a\}}.
\]
Descreva como isolar cada um dos finitos átomos $a$ de massa maior que
$\varepsilon$, usando pontos de continuidade $c<a<d$ suficientemente próximos
para que $F(a-)-F(c)<\varepsilon$ e $F(d)-F(a)<\varepsilon$; nos intervalos
restantes, refine a partição para que cada incremento de $F$ seja, por exemplo,
no máximo $2\varepsilon$. Use a monotonicidade e o controle nos pontos da
partição e nos limites à esquerda dos átomos isolados para obter a
convergência uniforme.
\end{enumerate}
\end{starredexercise}

\begin{starredexercise}\label{exercise:dois-a-dois}
Reveja a prova do \Cref{thm:lei-forte} e identifique exatamente em que ponto a
hipótese de independência dois a dois é usada. Mostre, por um exemplo, que a
prova do \Cref{thm:quarto-momento} não funciona com independência apenas dois a
dois, indicando qual passagem falha.
\end{starredexercise}

%----------------------------------------------------------------------
% BIBLIOGRAFIA
%----------------------------------------------------------------------
\bibliographystyle{plain}
\bibliography{../arquivos-aux/referencias}

\end{document}
